\chapter{Coherent comparison, descent and formal limits}
\label{cycmp:chapter}

Two fields can have the same product and different collision extensions.
Let $B=\mathbb C1\oplus\mathbb C\epsilon$, with $\epsilon^2=0$.
The constant fields
\[
 Y(a,u)b=ab,\qquad T=0,
\]
retain every ordered iterated product in $B$. They do not specify how
a particle in position three passes a particle in position two before
positions one and three coincide. Giving this passage either sign
produces the same binary coefficient and different ternary coefficients.
The distinction survives regular holonomicity, concentration in the
ordinary differential-module heart, and all local vacuum insertions.

We first construct these coefficients and identify the exchange datum
that makes their reconstruction fully faithful. We then turn to current
cochains, where a quadratic differential records multiplication in the
state algebra. Rational current fields require a jointly localized
tensor product, and their reconstruction must retain a current framing.
The final descent argument keeps the nonlinear comparison maps before
passing to finite cochains or to a formal limit.

\section{A collision invisible to the binary fields}
\label{cycmp:sec:particles}

Work over $k=\mathbb C$ on $X=\mathbb A^1$. Coefficients are enhanced
derived right algebraic differential modules with quasi-coherent
underlying modules. For a finite-type scheme $Z$, write
$J_Z=a_Z^!k$, where $a_Z:Z\to\operatorname{Spec}k$ is its structure map.
If $Z$ is smooth of dimension $n$, then $J_Z=\omega_Z[n]$.
We use cohomological shifts, so that $H^i(M[r])=H^{i+r}(M)$.

\begin{definition}\label{cycmp:def:vacua}
An ordered coefficient system with local vacua assigns $F_I$ to each
finite ordered set $I$, with $F_\varnothing=k$. A monotone map
$f:S\to I$ determines
\[
 d_f:X^I\longrightarrow X^S,
 \qquad (x_i)\longmapsto(x_{f(s)}),
 \qquad \nu_f:d_f^!F_S\longrightarrow F_I.
\]
These maps compose coherently and are equivalences when $f$ is
surjective. On the open set where every coordinate in $I_1$ differs
from every coordinate in $I_2$, a consecutive decomposition
$I=I_1\sqcup I_2$ has a specified isomorphism
\[
 F_{I_1}\boxtimes F_{I_2}\xrightarrow{\sim}F_I.
\]
It is compatible with refinements, monotone collisions and the maps
$\nu_f$.

For an active subset $S\subset I$, the vacuum-marked
coefficient is $p_{I,S}^!F_S$, where $p_{I,S}:X^I\to X^S$ is
projection. Its evaluation is $\nu_{S\hookrightarrow I}$.

The system satisfies the \emph{arity heart condition} when
$F_I[-|I|]$ is an ordinary right differential module. A morphism is a
degree-zero family preserving these operations and the identity of
the empty coefficient.
\end{definition}

Deleting an inactive position identifies two pullbacks of the same
active coefficient. It does not require $F_I$ itself to be pulled
back from a smaller configuration. This distinction permits a local
vacuum without imposing a scalar augmentation on its state algebra.
The shift in the heart condition is the canonical normalization of
\cite{BD04}*{\S2.1.1}.

Let $W$ be finite-dimensional and concentrated in degree zero. Choose
an involutive Yang--Baxter operator $c:W\otimes W\to W\otimes W$:
\[
 c^2=1,\qquad c_{12}c_{23}c_{12}=c_{23}c_{12}c_{23}.
\]
Here the subscripts specify the two tensor positions on which $c$
acts. Adjacent copies of $c$ define an action of the symmetric group
$S_r$ on $W^{\otimes r}$. Put $B=k1\oplus W$, with $W^2=0$, and
extend $c$ to $R_c:B^{\otimes2}\to B^{\otimes2}$ by
\[
 R_c(1\otimes b)=b\otimes1,
 \qquad R_c(b\otimes1)=1\otimes b,
 \qquad R_c|_{W^{\otimes2}}=c.
\]

For $0\le r\le n$, consider the reduced incidence scheme
\begin{equation}\label{cycmp:eq:incidence}
 Z_{n,r}=\left\{(x,t)\in X^n\times X^r:
  \prod_{i=1}^n(t_a-x_i)=0\ (1\le a\le r),\quad
  t_a\ne t_b\ (a\ne b)\right\}_{\mathrm{red}}.
\end{equation}
The selected roots $t_1,\ldots,t_r$ mark the occupied positions.
Unselected positions carry the vacuum. Write $p_{n,r}:Z_{n,r}\to X^n$
for projection, and set $Z_{n,0}=X^n$. The geometric permutation of
selected roots and the $c$-action on their states give an $S_r$-action.
Define
\begin{equation}\label{cycmp:eq:particles}
 F_n(W,c)=\bigoplus_{r=0}^n
       \left((p_{n,r})_*J_{Z_{n,r}}\otimes W^{\otimes r}\right)^{S_r}.
\end{equation}

\begin{theorem}\label{cycmp:thm:particles}
The coefficients \eqref{cycmp:eq:particles} form a translation-equivariant
ordered system with local vacua. Its consecutive factorization maps
are isomorphisms. Every coefficient is regular holonomic and satisfies
the arity heart condition. Its singleton is $J_X\otimes B$.
\end{theorem}

\begin{proof}
At a point of $Z_{n,r}$, let $I_a=\{i:x_i=t_a\}$. The sets $I_a$
are nonempty and disjoint. In a neighborhood of this point the
equations, after removing invertible factors, are
\[
 \prod_{i\in I_a}(x_i-t_a)=0\qquad(1\le a\le r).
\]

Use $x_i-t_a$, the $t_a$, and the remaining $x_i$ as coordinates.
The equations are square-free monomials in disjoint groups of
variables. They define a reduced local complete intersection of pure
dimension $n$. Thus $J_{Z_{n,r}}[-n]$ lies in the right-module heart.

Before the root inequalities, projection is finite because each root
satisfies a monic equation. Inverting the product of the root
differences makes $p_{n,r}$ affine and quasi-finite. Its direct image
is exact on that heart. These are the local complete-intersection and
direct-image statements of \cite{BD04}*{\S\S2.1.2,2.1.4}.
Finite sums and invariants are exact. Regular holonomicity follows
from preservation under the corresponding algebraic direct images
\cite{cycmp:HTT}*{Theorem 6.1.5}.

For a monotone surjection $f:S\twoheadrightarrow I$, reduced base
change of the incidence scheme replaces repeated factors in its
root polynomial by single factors. It is $Z_{|I|,r}$. Extraordinary
base change and pullback of the dualizing object therefore give
$d_f^!F_S\simeq F_I$. Differential modules are unchanged by the
discarded nilpotent thickenings.

For an active subset $S\subset I$, let $Y_{I,S,r}$ be the closed
reduced locus in $Z_{|I|,r}$ where every selected root occurs among
the coordinates of $S$. It is the reduced pullback of $Z_{|S|,r}$.
The closed embedding $h:Y_{I,S,r}\hookrightarrow Z_{|I|,r}$ has the
proper counit
\[
 h_*J_{Y_{I,S,r}}=h_*h^!J_{Z_{|I|,r}}\longrightarrow J_{Z_{|I|,r}}.
\]
Direct image and invariants give the vacuum evaluation.

Factor a
general monotone map through its image to obtain $\nu_f$. Inverse
images of active-root loci are the corresponding active-root loci.
Nested active subsets give nested closed embeddings. Base change
and transitivity of their proper counits prove composition, including
the empty active subset. The latter retains only $r=0$.

On a consecutive cross-disjoint open, each selected root belongs to
one block. Split the incidence scheme by that choice. Its pieces are
external products of the two incidence schemes, with $r=r_1+r_2$.
An invariant section on a free orbit of root orderings is determined
by its value at the ordering in which the first block comes first.
This identifies the $S_r$-invariants with the appropriate
$S_{r_1}\times S_{r_2}$-invariants, with coefficient one.
It gives the factorization isomorphism.

For several blocks, every parenthesization uses the same concatenated
root ordering. Alternative reorderings agree by the symmetric-group
relations. Splitting roots commutes with selecting active roots, so
the same proper base change proves compatibility with vacua and
collisions. These canonical enhanced-functor identities supply the
coherences for chains of maps. Simultaneous translation of all roots
and positions preserves every scheme and map used above. The empty
and singleton assertions follow directly from the formula.
\end{proof}

The normal extension makes the dependence on $c$ explicit. For a
surjection $q:V\to Q$ of finite-dimensional vector spaces, put
\[
 \mathcal P=k[u,u^{-1}]/k[u],\qquad
 M(q)=\ker\bigl(V[u,u^{-1}]\xrightarrow{q\otimes\mathrm{pp}}
                                      Q\otimes\mathcal P\bigr).
\]
Here $\mathrm{pp}$ takes the principal part. In the left-module
normal coordinate,
\begin{equation}\label{cycmp:eq:normal-lattice}
 M(q)=V[u]+(\ker q)\otimes u^{-1}k[u^{-1}],
 \qquad \frac vu\in M(q)\quad\Longleftrightarrow\quad q(v)=0.
\end{equation}

Multiplication by $u$ and differentiation preserve this module.
The right-module version is obtained by tensoring with the density.
Its supported quotient retains all normal derivatives, normalized by
$\operatorname{Res}(du/u)=1$ and
$u^{-r-1}=(-1)^r\partial_u^r(u^{-1})/r!$.
Completion replaces $k[u]$ by $k[[u]]$ and keeps finite Laurent poles.

\begin{lemma}\label{cycmp:lem:particle-lattices}
The binary normal module of $F_2(W,c)[-2]$ is $M(m_B)$.
Near $x_1=x_3\ne x_2$, with the output blocks ordered as $(13),2$,
the ternary normal module is $M(q_c)$, where
\begin{equation}\label{cycmp:eq:particle-outer}
 q_c=(m_B\otimes1)(1\otimes R_c):B^{\otimes3}\longrightarrow B^{\otimes2}.
\end{equation}
These descriptions include the collision quotient and the local vacua.
\end{lemma}

\begin{proof}
The union of the two sections $t=x_1$ and $t=x_2$ has, on the
normal line $u=x_1-x_2$, one additional principal-part quotient.
Each section has multiplicity one, and the quotient is
$(a(u),b(u))\mapsto\mathrm{pp}(a(u)+b(u))$. Its extra allowed pole
is $(1,-1)/u$. One can also obtain this description by dualizing
the exact sequence of constant sheaves on the union. The de Rham
comparison preserves that sequence and its maps
\cite{cycmp:HTT}*{Theorem 7.1.1}.

For two positions, the zero-particle summand is regular. The
one-particle summand is the preceding union of sections tensored
with $W$. The two-particle summand is fully localized, since two
selected roots cannot coincide. The quotient is consequently
\[
 1\otimes1\mapsto1,\quad w\otimes1\mapsto w,
 \quad1\otimes w\mapsto w,
 \quad w\otimes w'\mapsto0.
\]
This is $m_B$, independent of $c$. A vacuum includes one of the
two sections and is independent of $c$ as well.

At the outer collision, the middle root remains separate. The two
branches of a particle approaching the outer root have opposite
orders relative to a particle at the middle root. Expressing both
in the closed ordering $(13),2$ applies $R_c$ to positions two and
three. Multiplication then joins the outer positions. This is
$q_c$.

Terms with two particles at the outer positions remain fully
localized, and the other summands have the same two-section
description. The vacuum maps are the corresponding proper counits.
\end{proof}

Binary residue testing sends $f(u)a\otimes b$ to
$\mathrm{pp}(f(u)ab)$. Tests with negative powers recover the
regular coefficient $ab$ as well as every singular coefficient.
Thus the full binary fields are $Y(a,u)b=ab$, with $T=0$,
for every $c$. All planar iterates are products in $B$.

\begin{proposition}\label{cycmp:prop:particle-distinction}
For $W=k\epsilon$, the systems with $c=1$ and $c=-1$ have identical
binary coefficients, fields and local vacua. They are not isomorphic.
The identity of their binary fields does not lift to a coefficient map.
\end{proposition}

\begin{proof}
Set $v=\epsilon\otimes\epsilon\otimes1-
1\otimes\epsilon\otimes\epsilon$. Then $q_+(v)=0$, whereas
$q_-(v)=2\epsilon\otimes\epsilon$. Hence $v/(x_1-x_3)$ extends in
the first system and not in the second. Factorization would make an
identity unary lift the identity on the generic ternary tensor,
contradicting this extension test.

More generally a unital isomorphism of the state algebras sends
$\epsilon$ to $\lambda\epsilon$, with $\lambda\ne0$. Preserving the
outer quotient and its vacuum inclusions would require
$\lambda^2=-\lambda^2$. This is impossible in characteristic zero.
\end{proof}

\section{Reconstruction with an exchange}
\label{cycmp:sec:exchange}

The particle example suggests retaining the operator used to move a
state past a spectator. Multiplication can then be nonzero on two
non-vacuum states. No scalar augmentation is needed.

\begin{definition}\label{cycmp:def:exchange}
An algebra with involutive exchange is a nonzero finite-dimensional
unital associative algebra $B$, in degree zero, together with
$R:B^{\otimes2}\to B^{\otimes2}$ satisfying
\begin{align}
 R^2&=1,&R_{12}R_{23}R_{12}&=R_{23}R_{12}R_{23},\notag\\
 R(1\otimes a)&=a\otimes1,&R(a\otimes1)&=1\otimes a,
       &mR&=m,\label{cycmp:eq:exchange-unit}\\
 R(m\otimes1)&=(1\otimes m)R_{12}R_{23},&
 R(1\otimes m)&=(m\otimes1)R_{23}R_{12}.\label{cycmp:eq:exchange-fusion}
\end{align}
Morphisms are unital algebra maps intertwining $R$.
\end{definition}

These are identities on state tensors. A tensor category of
representations, with its own exchange, remains additional structure.
For a function $f:I\to J$ of finite ordered sets, sort the inputs
into increasing $f$-fibers using adjacent copies of $R$.
Multiply inside each fiber and insert $1$ in empty fibers.
Denote this map by
\[
 Q_f:B^{\otimes I}\longrightarrow B^{\otimes J}.
\]

\begin{lemma}\label{cycmp:lem:fusion}
The maps $Q_f$ satisfy $Q_gQ_f=Q_{gf}$ and $Q_{\mathrm{id}}=1$.
They preserve consecutive external products. If $f$ is surjective,
$Q_f$ has the section $Q_s$ for any set-theoretic section $s$ of $f$.
\end{lemma}

\begin{proof}
Involution, Yang--Baxter and commutation in disjoint slots give a
symmetric-group action. An exchange inside a single fiber disappears
after multiplication, by associativity and $mR=m$.

In $Q_gQ_f$, use \eqref{cycmp:eq:exchange-fusion} to move each
multiplication through the exchanges above it. Each move replaces a
crossing of products by crossings of their separate inputs. The
process terminates at individual inputs.

The unit equations remove
empty fibers, leaving the permutation that sorts by $gf$ and the
products in its fibers. Associativity removes parentheses, and
$mR=m$ removes permutations internal to a final fiber. This is
$Q_{gf}$. Juxtaposition gives the external-product identity, and
$fs=\mathrm{id}$ gives $Q_fQ_s=1$. Longer composites have the same
map $Q$, so the same argument supplies their coherences.
\end{proof}

Every unital associative algebra admits an example:
\begin{equation}\label{cycmp:eq:context-exchange}
 R_B(a\otimes b)=ab\otimes1+1\otimes ab-a\otimes b.
\end{equation}
Indeed, for $p(a\otimes b)=ab\otimes1$ and
$q(a\otimes b)=1\otimes ab$, one has
$p^2=p$, $q^2=q$, $pq=p$ and $qp=q$. Hence $(p+q-1)^2=1$.
Both braid words on $a\otimes b\otimes c$ give
\[
\begin{split}
 &abc\otimes1\otimes1-ab\otimes c\otimes1
       +1\otimes1\otimes abc+1\otimes ab\otimes c\\
 &\qquad+a\otimes bc\otimes1-1\otimes a\otimes bc
       -a\otimes b\otimes c.
\end{split}
\]

The first fusion identity gives
$abc\otimes1+1\otimes abc-ab\otimes c$ on either side.
The second gives $abc\otimes1+1\otimes abc-a\otimes bc$.
The remaining identities follow from the unit and associativity.

We now specify the collision extension independently of its fields.
Let $X$ be a connected smooth separated complex variety of positive
dimension. An equality partition $\pi$ of $[n]$ specifies the stratum
$S_\pi\subset X^n$ on which coordinates agree exactly within its
blocks. Order those blocks by first occurrence, let
$f_\pi:[n]\to[|\pi|]$ be the quotient, and put $q_\pi=Q_{f_\pi}$.
For a coarsening $\rho$ of $\pi$, fusion on their block sets gives
\[
 q_\rho=q_{\pi\rho}q_\pi,\qquad
 q_{\rho\sigma}q_{\pi\rho}=q_{\pi\sigma}.
\]

Write $j_n:\operatorname{Conf}_n(X)\hookrightarrow X^n$ for the
fully distinct locus and set
\[
 E_n=(B^*)^{\otimes n},\qquad E_\pi=\operatorname{im}(q_\pi^*)\subset E_n.
\]
Coarsening makes the subspace smaller.

Define $K_R(B)_n$ as the subsheaf of the ordinary analytic sheaf
$(j_n)_*\underline E_n$ whose germ on $S_\pi$ belongs to $E_\pi$.
Near such a point choose disjoint coordinate neighborhoods of the
distinct positions. The fully distinct locus in their product is
connected, since complex diagonals have real codimension at least
two. Its generic germ is a single vector.

Every vector in $E_\pi$
extends constantly nearby: the nearby partitions refine $\pi$ and
their subspaces contain $E_\pi$. Thus the stalk is exactly $E_\pi$,
and its generizations are the specified inclusions. This is an
algebraically constructible sheaf.

Normalize the regular Riemann--Hilbert equivalence by
$\operatorname{DR}(\omega_Y)=k_{Y^{\mathrm{an}}}[\dim Y]$ on a smooth
variety. Let $\mathbb D$ denote Verdier duality, and define
\begin{equation}\label{cycmp:eq:exchange-coefficient}
 F_R(B)_n=\operatorname{DR}^{-1}(\mathbb D K_R(B)_n),
 \qquad F_R(B)_0=k.
\end{equation}
Duality and the six operations are compatible with this equivalence
\cite{cycmp:HTT}*{Theorems 7.1.1,7.2.2}.

\begin{theorem}\label{cycmp:thm:exchange-reconstruction}
Formula \eqref{cycmp:eq:exchange-coefficient} gives a fully faithful
functor from finite-dimensional algebras with involutive exchange to
ordered coefficient systems with local vacua on $X$.
Its coefficients are bounded regular-holonomic complexes, and its
consecutive factorization maps are isomorphisms. For
$\iota_\pi:S_\pi\hookrightarrow X^n$,
\[
 \iota_\pi^!F_R(B)_n\simeq J_{S_\pi}\otimes B^{\otimes|\pi|}.
\]
On the affine line it is translation-equivariant, has unary translation
zero, and satisfies the arity heart condition.
\end{theorem}

\begin{proof}
The constructible sheaf has finite-dimensional stalks, so its Verdier
dual corresponds to a bounded regular-holonomic complex. Duality
turns stalk restriction into extraordinary restriction. Since
$q_\pi$ is split surjective, the dual of $E_\pi$ is canonically
$B^{\otimes|\pi|}$, proving the displayed formula.

On a monotone collision diagonal, equality partitions are the inverse
images of partitions of the smaller ordered set. Their first-occurrence
orders agree. The sheaf stalks and generizations therefore identify
the ordinary pullback with the smaller sheaf. Duality gives the
collision equivalence.

For an active subset, use the injection of
its equality-block set into the full block set. The map $Q$ for this
injection moves the active factors into the required order and inserts
units elsewhere. Its dual gives
$K_R(B)_I\to p_{I,S}^*K_R(B)_S$. After duality this is the local
vacuum evaluation. Fusion and insertion commute because both
composites are $Q$ of the same function.

On a consecutive cross-disjoint open, every equality block lies on
one side of the cut. The quotient and its dual image are external
products of the two block constructions. This gives strong
factorization.

Refinements, collisions and vacuum insertions again
reduce to equalities of $Q$-maps. Enhanced derived functors carry
these sheaf identities to coherent structure maps. A morphism of
exchange algebras intertwines all $Q_f$. Its dual tensor powers and
then duality give the claimed covariant functor.

For heart concentration on a curve, filter $B$ by
$L_0B=k1$, $L_rB=B$ for $r\ge1$, and tensor powers by total
filtration. Each $Q_f$ preserves the filtration. For surjective $f$
its filtered section proves strictness. The associated graded is
$k1\oplus W$, $W=B/k1$, with $W^2=0$. Its exchange has the form
$R_c$ of the particle construction.

The decreasing dual filtration
on $E_n$ induces on each $E_\pi$ the dual output filtration.
Strictness identifies its associated graded inclusion with
$(\operatorname{gr}q_\pi)^*$.

The graded sheaf is the Verdier dual of
\eqref{cycmp:eq:particles}. In particle degree $r$, its stalk selects
$r$ distinct equality blocks and one $W^*$ factor at each. Under
generization these blocks select finer labels, reordered by $c^*$.
These are exactly the components of the graded inclusion. The
compactly supported direct image from the selected-root cover has
finite fibers, so this identification has no higher stalk terms.
Evaluation at one ordering gives coefficient one.

The particle
theorem puts every dual graded piece in the shifted right heart.
There are at most $n+1$ filtration steps. Closure of the heart under
extensions proves concentration for $F_R(B)_n[-n]$.

It remains to identify morphisms. The singleton is constant, and
$X$ is connected, so a unary map is a constant linear map $f:B\to B'$.
The vacuum makes it unital, and factorization makes its generic
arity-$n$ map $f^{\otimes n}$. Binary generization imposes
$fm=m'(f\otimes f)$.

At $13\mid2$ the quotient is
\[
 q_R=(m\otimes1)(1\otimes R),\qquad
 q_R(a\otimes b\otimes1)=a\otimes b,
 \quad q_R(1\otimes a\otimes b)=R(a\otimes b).
\]
If $h$ is the closed coefficient map, naturality gives
$hq_R=q_{R'}f^{\otimes3}$. The first vacuum insertion forces
$h=f\otimes f$, and the second forces
$(f\otimes f)R=R'(f\otimes f)$.

Conversely, such an $f$ has the structured lift already constructed.
After duality, the difference between two lifts would have zero
generic restriction and image in a sheaf embedded in its generic
direct image. It is zero.

The sheaves and the domains of their
structure maps, formed by ordinary pullbacks and external products,
are in degree zero. Their mapping spaces are discrete, and limits
imposing compatibility remain discrete. This proves full faithfulness
with its coherent comparisons.

Simultaneous translation on the line
preserves every equality stratum and subspace. The singleton has
constant connection, giving $T=0$.
\end{proof}

This theorem recovers an exchange from a vacuum-marked outer
attachment. It does not assert that the constant binary fields
determine that attachment. Matrix states make the difference visible
without a square-zero multiplication.

\begin{example}\label{cycmp:ex:matrix-exchanges}
In $B=\operatorname{Mat}_2(k)$, let
\[
 X_0=\begin{pmatrix}0&1\\1&0\end{pmatrix},\qquad
 Z_0=\begin{pmatrix}1&0\\0&-1\end{pmatrix},\qquad
 e_{a,b}=X_0^aZ_0^b,
\]
where $(a,b)\in(\mathbb Z/2)^2$. Multiplication is
$e_{a,b}e_{c,d}=(-1)^{bc}e_{a+c,b+d}$.
Besides \eqref{cycmp:eq:context-exchange}, define
\[
 R_{\mathrm{col}}(e_{a,b}\otimes e_{c,d})
       =(-1)^{bc+ad}e_{c,d}\otimes e_{a,b}.
\]

The exponent is symmetric and bilinear. It proves involution, both
fusion identities and Yang--Baxter. The latter has on either side
the sum of the three pairwise exponents. Multiplication after this
exchange has sign $(-1)^{bc+ad+da}=(-1)^{bc}$, as required.
All states remain in cohomological degree zero.

Both extensions have binary quotient $m$ and fields $Y(a,u)b=ab$.
Let $s(a\otimes b)=a\otimes b\otimes1$ and put
$v=1\otimes X_0\otimes Z_0-sR_B(X_0\otimes Z_0)$.
The outer quotient for $R_B$ kills $v$, while that for
$R_{\mathrm{col}}$ gives
\[
 -Z_0\otimes X_0-X_0Z_0\otimes1-1\otimes X_0Z_0+X_0\otimes Z_0.
\]
Its four terms are distinct basis tensors. Hence the two normal
modules admit different poles $v/(x_1-x_3)$.

They are not related by a change of unary framing either:
$\operatorname{tr}R_B=4+4-16=-8$, whereas
$\operatorname{tr}R_{\mathrm{col}}=4$. An isomorphism would conjugate
one exchange to the other by $f\otimes f$.

For $P=E_{12}$ and $Q=E_{21}$, both exchanges retain the ordered
outer product $PQP=P$. An ordinary flip would give $P^2Q=0$.
There is also no unital algebra map $B\to k$: it would kill $P,Q$
by nilpotence and contradict $PQ+QP=1$. Thus the local vacua of
the coefficient system cannot be replaced by a scalar counit on $B$.
\end{example}

\section{Multiplication recorded by a current differential}
\label{cycmp:sec:currents}

The preceding reconstruction retains exchange as part of the input.
A different construction records a commutative multiplication inside
a quadratic differential. Its scalar object already explains why the
empty field unit and the coefficient of constant states must be kept
separate.

Let $\mathcal A$ consist of unital commutative dg algebras in
nonpositive degrees, with a degree-zero derivation $T_A$ commuting
with $d_A$. Require
\[
 A=\bigoplus_{q\ge0}A_q,\qquad A_pA_q\subset A_{p+q},
 \quad d_AA_q\subset A_q,\quad T_AA_q\subset A_{q+1},
\]
where every $A_q$ is a bounded finite-dimensional complex and
$1\in A_0$. Maps preserve the unit, differential and translation.
They need not preserve the auxiliary weight decomposition.
The direct product $A\times A'$ has componentwise multiplication
and underlying complex $A\oplus A'$. We admit the zero algebra as
its monoidal unit.

Fix $\mathfrak g=\mathfrak{sl}_2(k)$, with basis $e,f,h$ and
$[h,e]=2e$, $[h,f]=-2f$, $[e,f]=h$. Define
\begin{equation}\label{cycmp:eq:current-lie}
 L_A=\mathfrak g\otimes A,\quad
 [x\otimes a,y\otimes b]=[x,y]\otimes ab,
 \quad d_L=1\otimes d_A,\quad T_L=1\otimes T_A.
\end{equation}
Its Chevalley cochain algebra is
\begin{equation}\label{cycmp:eq:current-cochains}
 \mathcal B_A=\widehat{\operatorname{Sym}}(W_A),
 \qquad W_A=(L_A[1])^\vee.
\end{equation}

On a discrete graded vector space, $\vee$ means its full graded
dual with the product topology. On a linearly compact space it means
the continuous dual. The completion is the product over weights
in each cohomological degree. A generator dual to $L_A^i$ has
degree $1-i>0$, so only finitely many symmetric degrees occur in
each cohomological degree.

We fix signs before using the differential. For a homogeneous basis
$v_i$ of $L_A$, put $a_i=|v_i|$ and write
\[
 d_Lv_i=\sum_k D_i^k v_k,\qquad
 [v_i,v_j]=\sum_k f_{ij}^k v_k.
\]
The dual generator $c^i$ has degree $1-a_i$. The cochain
differential is the continuous derivation
\begin{equation}\label{cycmp:eq:ce-signs}
 dc^k=\sum_i(-1)^{a_i}D_i^k c^i
       -\frac12\sum_{i,j}(-1)^{a_i(1-a_j)}f_{ij}^k c^ic^j.
\end{equation}

Its linear and quadratic parts are denoted $d_1,d_2$.
The dg Lie identities give
$d_1^2=d_1d_2+d_2d_1=d_2^2=0$.
For degree-zero currents this says
$(d\lambda)(x,y)=-\lambda([x,y])$.

Set $T_{\mathcal B}=-T_L^\vee$ on generators and extend it as a
derivation. The regular fields are
\begin{equation}\label{cycmp:eq:regular-fields}
 Y_{\mathcal B}(b,u)c=(e^{uT_{\mathcal B}}b)c.
\end{equation}

The completed recipient is explicit: an output through dual weight
$-q$ and Taylor degree $N-1$ needs inputs only through weight
$q+N-1$. Equation \eqref{cycmp:eq:regular-fields} denotes the
compatible operations on these finite quotients. It does not
impose a uniform bound on the input weights. Since $T_L$ is a dg
Lie derivation, $T_{\mathcal B}$ is a product derivation commuting
with $d$. Associativity, graded commutativity and the exponential
identity give the dg vertex-superalgebra identities. For example,
weak associativity is
\[
 (e^{(u+v)T_{\mathcal B}}b)(e^{vT_{\mathcal B}}c)a
 =e^{vT_{\mathcal B}}\bigl((e^{uT_{\mathcal B}}b)c\bigr)a.
\]

Constant currents give $\rho_A:\mathcal B_A\to\mathcal B_k$.
The latter algebra is
\[
 \mathcal B_k=\Lambda(\xi_e,\xi_f,\xi_h),\qquad
 d\xi_e=2\xi_e\xi_h,\quad
 d\xi_f=-2\xi_f\xi_h,\quad d\xi_h=-\xi_e\xi_f,
 \qquad T_{\mathcal B_k}=0.
\]
The differential identifies degrees one and two and vanishes in
degrees two and three. Thus $H^*(\mathcal B_k)$ is $k$ in
degrees zero and three. It is different from the field unit
$k=\mathcal B_0$. An augmentation $A\to k$ commuting with
translation gives a section of $\rho_A$.

\begin{definition}\label{cycmp:def:current-presentation}
A quadratic current presentation is a complete regular dg field
algebra $\mathcal B=\widehat{\operatorname{Sym}}(W)$ with a specified
split symmetric-degree grading and a map
$\rho:\mathcal B\to\mathcal B_k$. Its generators have the form
$W=(L[1])^\vee$, with $L$ nonpositive and with bounded finite
nonnegative weight pieces. The differential has only linear and
quadratic parts, with the signs \eqref{cycmp:eq:ce-signs}.
Translation preserves $W$ and has the stated weight behavior.
The fields are \eqref{cycmp:eq:regular-fields}.

The transpose of $\rho$ is a dg Lie map $j:\mathfrak g\to L$,
with $T_Lj=0$. Require evaluation to be an isomorphism
\begin{equation}\label{cycmp:eq:isotypic}
 \mathfrak g\otimes\operatorname{Hom}_{\mathfrak g}(\mathfrak g,L)
       \xrightarrow{\sim}L,
 \qquad x\cdot v=[j(x),v].
\end{equation}
The recovered multiplicity space has the weight properties of
$\mathcal A$. The map $\rho$ preserves weight and symmetric degree.
Morphisms are continuous dg field maps over $\mathcal B_k$ that
preserve the split symmetric-degree grading. Denote this category
by $\mathcal P_{\mathfrak g}$.
\end{definition}

The split grading is supplied data. It makes morphisms linear on
generators even when a generator has cohomological degree greater
than one. No assertion that every dg field algebra has such a
canonical grading is needed.

\begin{theorem}\label{cycmp:thm:current-recognition}
There is an equivalence
\[
 \mathcal A^{\mathrm{op}}\xrightarrow{\sim}\mathcal P_{\mathfrak g},
 \qquad A\longmapsto(\mathcal B_A,\rho_A).
\]
Its inverse recovers
\[
 A_{\mathcal B}=
 \operatorname{Hom}_{\mathfrak g}(\mathfrak g,W^\vee[-1]).
\]
The linear differential recovers $d_A$, the quadratic differential
recovers multiplication, and $-T_{\mathcal B}^\vee$ recovers $T_A$.
Direct products correspond to completed tensor products, whose
scalar map uses multiplication in $\mathcal B_k$.
\end{theorem}

\begin{proof}
The adjoint representation of $\mathfrak{sl}_2$ satisfies
\[
 \operatorname{End}_{\mathfrak g}(\mathfrak g)=k,
 \qquad
 \operatorname{Hom}_{\mathfrak g}(\mathfrak g\otimes\mathfrak g,
                                \mathfrak g)=k[\,,\,].
\]
The first identity follows from the weight spaces $2,0,-2$ and
the actions of $e,f$ between them. For the second,
$\Lambda^2\mathfrak g\simeq\mathfrak g$, whereas
$\operatorname{Sym}^2\mathfrak g$ is the five-dimensional
irreducible representation generated by $e\otimes e$ plus
the invariant line generated by
$2(e\otimes f+f\otimes e)+h\otimes h$.

Put $M=\operatorname{Hom}_{\mathfrak g}(\mathfrak g,L)$.
The bracket on $\mathfrak g\otimes M$ therefore has the form
$[x,y]\otimes m(a,b)$. Graded antisymmetry makes $m$ graded
commutative. Jacobi on $e\otimes a,e\otimes b,f\otimes c$ gives
$a(bc)=(-1)^{|a||b|}b(ac)$. Consequently
\[
 (ab)c=(-1)^{(|a|+|b|)|c|}c(ab)
      =(-1)^{|b||c|}a(cb)=a(bc).
\]
The prescribed adjoint action makes $j\in M^0$ the unit.

Both $d_L$ and $T_L$ commute with that action. The first adjoint
identity makes them $1\otimes d_M$ and $1\otimes T_M$.
Their Lie derivation equations are precisely the product
derivation equations on $M$.

The symmetric-degree components of $d^2=0$ recover the differential,
derivation and Jacobi equations of $L$. A morphism of presentations
transposes on generators to a dg Lie map preserving $j$ and
translation. It is $1\otimes f$, and its bracket equation says
exactly that $f$ preserves multiplication. Continuous duality and
the free completed presentation recover the whole field map from
this linear map. This proves the inverse comparisons and full
faithfulness.

Finally $L_{A\times A'}=L_A\oplus L_{A'}$ gives
\[
 \mathcal B_{A\times A'}\simeq
       \mathcal B_A\widehat\otimes\mathcal B_{A'},\qquad
 \mathcal B_0=k.
\]
The diagonal constant-current inclusion transposes to multiplication
of the two scalar maps, as asserted.
\end{proof}

\section{The full coefficient comparison}
\label{cycmp:sec:regular-comparison}

Recovering the state multiplication is only the first step. We must
also recover the normal derivatives and compare all ways of reaching
a multiple collision. The elementary three-point identity fixes the
recipient of that comparison.

Put $s=z_1-z_3$, $t=z_2-z_3$ and $r=s-t$. For a vector space $P$,
let
\[
 \mathcal G_3(P)=P[[s,t]][(st(s-t))^{-1}],
 \qquad E_{s,t}(P)=\bigoplus_{a,b\ge1}P\,s^{-a}t^{-b}.
\]
The latter is the full normal contact module, with multiplication by
coordinates followed by principal part and with normal differentiation.
Change from contact coordinates $(r,t)$ to $(s,t)$ by
\begin{equation}\label{cycmp:eq:contact-change}
 J(r^{-a}t^{-b})=
   \sum_{j=0}^{b-1}\binom{a+j-1}{j}s^{-a-j}t^{-b+j}.
\end{equation}
This map sends the simple contact to $s^{-1}t^{-1}$, carries
$\partial_r$ to $\partial_s$, and carries the old $\partial_t$
to $\partial_s+\partial_t$. Applying these derivatives to the
simple contact proves the formula for every $a,b$.

\begin{lemma}\label{cycmp:lem:three-channels}
For $F\in\mathcal G_3(P)$, define
\[
\begin{aligned}
 L(F)&=J\,\mathrm{pp}_t\mathrm{pp}_rF(r+t,t)
                                     &&(r\ll t),\\
 R(F)&=\mathrm{pp}_s\mathrm{pp}_tF(s,t)
                                     &&(t\ll s),\\
 N(F)&=\mathrm{pp}_t\mathrm{pp}_sF(s,t)
                                     &&(s\ll t).
\end{aligned}
\]
The rightmost principal part is applied first. Then
$L(F)-R(F)+N(F)=0$ in $E_{s,t}(P)$.
\end{lemma}

\begin{proof}
Write $F=G/(s^at^b(s-t)^c)$. A monomial surviving two
principal parts has total degree at most $-2$. The denominator
is homogeneous, so numerator terms above degree $a+b+c-2$
contribute to none of the three maps. We may take $G$ polynomial.
Over $k[t,t^{-1}]$, partial fractions in the coprime factors
$s,s-t$ reduce to polynomials in $s$, terms $a_j(t)s^{-j}$,
and terms $b_j(t)(s-t)^{-j}$, with Laurent coefficients in $t$.

A polynomial in $s$ contributes nothing. For $a_j(t)s^{-j}$,
the $r$-expansion is regular, while both other channels give
$s^{-j}\mathrm{pp}_t(a_j)$.

For $b_j(t)(s-t)^{-j}$, the
expansion for $s\ll t$ is regular in $s$. In the other domain,
\[
 (s-t)^{-j}=\sum_{n\ge0}\binom{j+n-1}{n}s^{-j-n}t^n.
\]
A coefficient $t^{-\beta}$ contributes only for $\beta\ge1$,
and gives
$\sum_{n<\beta}\binom{j+n-1}{n}s^{-j-n}t^{n-\beta}$.
This is exactly $J(r^{-j}t^{-\beta})$. The identity follows
term by term, with the coefficients of $P$ unchanged.
\end{proof}

For $A\in\mathcal A$, put $\mathscr M_A=\omega_X\otimes A$,
with its internal differential and right action
\begin{equation}\label{cycmp:eq:right-current}
 (f\,dz\otimes a)\partial_z
       =-f'\,dz\otimes a-f\,dz\otimes T_Aa.
\end{equation}
On an ordered chart, a partition $\pi$ of $I$ has its blocks
in root order. Let $i_\pi$ be its diagonal inclusion and $j_\pi$
the open inclusion of distinct block positions. Define the
partition coefficient complex
\begin{equation}\label{cycmp:eq:regular-coefficient}
 (C_A)_I=
 \bigoplus_\pi i_{\pi,*}j_{\pi,*}j_\pi^!
                  (\mathscr M_A^{\boxtimes|\pi|})[|\pi|],
 \qquad (C_A)_\varnothing=k.
\end{equation}
All direct images and pullbacks are derived. In particular the closed
direct image includes every normal derivative.

Regard each block state as suspended, with $|sa|=|a|-1$ and
$d(sa)=-s(d_Aa)$. A merger $i<j$ moves $sa_i$ beside $sa_j$
through the intervening suspended factors, then applies
\begin{equation}\label{cycmp:eq:regular-residue}
 sa\otimes sb\longmapsto
    (-1)^{|a|}s\,\mathrm{pp}_u\bigl(f(u)(e^{uT_A}a)b\bigr),
 \qquad u=z_i-z_j,
\end{equation}
retaining the root of $j$. The sign of tensoring this degree-one
map past the preceding factors is included. Equivalently, the total
sign is
\[
 (-1)^{\sum_{k<i}(|a_k|-1)+
       |a_i|(1+\sum_{i<k<j}(|a_k|-1))}.
\]
It is $(-1)^{i-1}$ for degree-zero states. The differential of
$C_A$ is the internal differential plus these mergers.

Use the full ordered collision carrier of
Chapter~\ref{cycuts:ch}: all ordered-subset cuts are retained,
including $13\mid2$. Denote its all-subset convolution by $\star$
and its empty unit by $K$. Write $\sharp$ for the full coefficient
duality of Theorem~\ref{cyjets:thm:pro}, and $\star_\sharp$ for
the tensor transported by that duality. These are differential-module
duals. They are distinct from the continuous vector-space dual in
\eqref{cycmp:eq:current-cochains}.

\begin{proposition}\label{cycmp:prop:regular-coefficients}
The complexes $C_A$ form a functor to full ordered collision
coalgebras with local vacua. It sends quasi-isomorphisms in
$\mathcal A$ to equivalences. There are coherent isomorphisms
\begin{equation}\label{cycmp:eq:colored-coefficients}
 C_{A\times A'}\simeq C_A\star C_{A'},\qquad C_0=K.
\end{equation}
Under this comparison the diagonal $A\to A\times A$ induces
the full cut coproduct of $C_A$.
\end{proposition}

\begin{proof}
The regular fields of $A$ are $Y_A(a,u)b=(e^{uT_A}a)b$.
For three states their iterates are expansions of the same
graded-commutative product
$(e^{sT_A}a)(e^{tT_A}b)c$. Multiplication by an allowed rational
test puts them in the common recipient of
Lemma~\ref{cycmp:lem:three-channels}. That lemma, with the other
positions as parameters, gives the three-channel residue identity.

Two disjoint mergers cancel in $d^2$ because their suspended maps
have degree one. The three mergers on a fixed triple cancel by
the displayed identity and the same suspended permutations. These
exhaust pairs of mergers. The internal differential anticommutes
with them since $d_A$ is a product derivation commuting with $T_A$.
Thus the proposed differential squares to zero.

Cuts, diagonal restriction and exchange of ordered charts retain
these same residue maps and their full normal-transfer modules.
The unit identities follow from $Y_A(1,u)a=a$ and
$Y_A(a,u)1=e^{uT_A}a$. Associativity and the three-channel identity
give compatible compositions on successive collisions. Ordinary
graded exchanges give the chart coherences. This constructs the
full coefficient coalgebra and its local vacua.

Filter a coefficient with $n$ labels by its number of blocks.
The internal differential preserves this number and a merger lowers
it. There are only $n$ nonempty steps.

A quasi-isomorphism $A\to A'$
gives an equivalence of the singleton differential modules.
Finite external products and derived geometric functors preserve
equivalences, giving the same conclusion on every associated graded term.
Finite filtered comparison gives the assertion on every configuration. The natural
collision, cut and vacuum maps are unchanged by this argument.

For \eqref{cycmp:eq:colored-coefficients}, expand every block state
of $A\oplus A'$ in its two colors. The unions $I_0,I_1$ of the
blocks of each color specify exactly an all-subset cut. The two
internal distinct-block opens, together with every cross-color
diagonal inverted by convolution, give the original open. The
closed direct images and normal modules agree.

Group suspended factors by color with their Koszul permutation sign.
Opposite colors have zero product. Same-color mergers and internal
differentials are those of the corresponding factor. Moving a
degree-one merger through the other color gives precisely the tensor
differential sign. Thus grouping is a chain isomorphism.

Collisions
either identify same-color blocks or kill a cross-color localized
term. Cuts keep the same colored blocks. Grouping three colors in
either order crosses each pair once, proving associativity, symmetry
and the unit law.

Finally the diagonal expands each block in both
colors. A given cut keeps the term when every block stays on one
side and gives zero otherwise. These are the full coproduct maps.
\end{proof}

Define $E_A=C_A^\sharp$ and, for a quadratic current presentation,
\[
 \mathfrak F(\mathcal B)=C_{A_{\mathcal B}}^\sharp.
\]
Recognition reverses arrows and full coefficient duality reverses
them again.

\begin{theorem}\label{cycmp:thm:regular-realization}
The functor $\mathfrak F$ is covariant and strong symmetric monoidal.
There is a natural isomorphism $\mathfrak F(\mathcal B_A)\simeq E_A$
preserving every collision, all normal derivatives, arbitrary spectators
and every ordered-subset cut. It carries field multiplication to
$\Delta_A^\sharp$ in both orders. The field unit maps to $K$, while
\[
 \mathfrak F(\mathcal B_k)=C_k^\sharp,
 \qquad \mathfrak F(\rho_A):E_A\longrightarrow C_k^\sharp.
\]
Any scalar section induced by an augmentation of $A$ is also preserved.
\end{theorem}

\begin{proof}
Recognition identifies the recovered state algebra with $A$, including
its differential and translation. Apply that identification to each
partition block and each normal derivative in
\eqref{cycmp:eq:regular-coefficient}. It is a coefficient chain
isomorphism before duality, compatible with all structure maps.
The tensor comparison is the transpose of
\eqref{cycmp:eq:colored-coefficients}, together with the cochain
tensor comparison. Their coherence identities were proved above.

On generators, multiplication of cochains transposes to the diagonal
$L_A\to L_A\oplus L_A$. Its recovered state map is
$A\to A\times A$, so the last assertion of the proposition identifies
the resulting product with $\Delta_A^\sharp$. Color interchange
proves the opposite-order assertion. The zero algebra and the maps
$k\to A\to k$, when the latter is supplied, give the stated empty
and local scalar maps.

The translation sign is fixed: $e^{uT_{\mathcal B}}$ transposes to
$e^{-uT_A}$. Its coefficient map is $C(e^{-uT_A})^\sharp$ at every
finite Taylor precision. For a normal test $u^{-m}$, the complete
residue is the finite expression
\begin{equation}\label{cycmp:eq:all-normal-regular}
 r(u^{-m};a,b)=\sum_{j=0}^{m-1}\frac{(T_A^ja)b}{j!}u^{j-m}.
\end{equation}
This compares every normal derivative, including at an outer collision
with spectators. The singleton is $(\mathscr M_A[1])^\sharp$.
It is obtained from the recovered linear generators, not from one
uniform shift of the entire Chevalley cochain algebra.
\end{proof}

\section{Formal equivalence and the linearized states}
\label{cycmp:sec:formal}

The current differential recovers the state algebra before cohomology
is taken. To make reconstruction invariant under a change of dg
presentation, its weak equivalences must retain the linearized states.
Nilpotent test algebras detect exactly this requirement.

Let $\mathfrak a$ be the augmentation ideal of a quadratic current
presentation, and put
\[
 F^r\mathcal B=\overline{\mathfrak a^r},\qquad
 Q\mathcal B=\mathfrak a/\overline{\mathfrak a^2}.
\]
The differential on $Q\mathcal B=W$ is $d_1$.

Let $R$ be a finite-dimensional local commutative dg algebra in
nonpositive degrees, augmented to $k$, with nilpotent augmentation
ideal $\mathfrak m$. The polynomial de Rham algebra of the
$n$-simplex is
\[
 \Omega_n=k[t_0,\ldots,t_n,dt_0,\ldots,dt_n]
      /(\textstyle\sum t_i-1,\sum dt_i),
 \qquad |t_i|=0,\quad |dt_i|=1.
\]

Define a simplicial set of formal points by
\begin{equation}\label{cycmp:eq:formal-points}
 \mathcal X_{\mathcal B}(R)_n=
 \left\{\phi:\mathcal B\longrightarrow R\otimes\Omega_n:
 \begin{array}{l}
 \phi\text{ is a continuous dg algebra map},\\
 \phi(\mathfrak a)\subset\mathfrak m\otimes\Omega_n
 \end{array}\right\}.
\end{equation}
The target is discrete. Faces and degeneracies are induced by the
simplex maps. The scalar map gives
$\mathcal X_{\mathcal B_k}\to\mathcal X_{\mathcal B}$.

A morphism $f:\mathcal B\to\mathcal B'$ in
$\mathcal P_{\mathfrak g}$ is a \emph{formal equivalence} when
$\mathcal X_{\mathcal B'}(R)\to\mathcal X_{\mathcal B}(R)$ is a
weak equivalence for every such $R$. This definition uses the field
algebra and its augmentation. It does not use the coefficient functor.

\begin{theorem}\label{cycmp:thm:formal-equivalence}
For a morphism of quadratic current presentations, the following
conditions are equivalent:
\begin{enumerate}[label=\textup{(\roman*)}]
\item it is a formal equivalence.
\item its map on $Q\mathcal B$ is a quasi-isomorphism.
\item its recovered map of state algebras is a quasi-isomorphism.
\end{enumerate}
Each condition implies an ordinary quasi-isomorphism of field algebras.
\end{theorem}

\begin{proof}
The images of the generators under a map in
\eqref{cycmp:eq:formal-points} determine a total degree-one element
$\alpha\in L\otimes\mathfrak m\otimes\Omega_n$.
The signs in \eqref{cycmp:eq:ce-signs} identify the chain-map equation
with $d\alpha+\tfrac12[\alpha,\alpha]=0$.

Continuity means finite weight support on each simplex. Only finitely
many generator degrees occur, since $R\otimes\Omega_n$ has bounded
degree range. Conversely, a Maurer--Cartan element determines a
continuous map on the generators. Nilpotence of $\mathfrak m$ makes
its extension to the completed symmetric algebra well defined.
Thus
\[
 \mathcal X_{\mathcal B}(R)=
 \operatorname{MC}_\bullet(L\otimes\mathfrak m).
\]

The dg Lie algebra on the right is nilpotent and nonpositive.
A quasi-isomorphism of $L$ therefore induces a weak equivalence
of these simplicial sets by \cite{cycmp:Getzler}*{Theorem 4.8}.

Conversely take $R=k[\varepsilon]/(\varepsilon^2)$, in degree zero.
The bracket on $L\varepsilon$ vanishes. Polynomial closed forms
give the Dold--Kan realization of the shifted complex, so
\begin{equation}\label{cycmp:eq:tangent-homotopy}
 \pi_n\mathcal X_{\mathcal B}(R)
       =H^{1-n}(L)\varepsilon\quad(n\ge1),\qquad \pi_0=*.
\end{equation}
See \cite{cycmp:Hinich}*{\S9.2} for this polynomial-forms comparison.
Every degree of $L$ is nonpositive and hence occurs in
\eqref{cycmp:eq:tangent-homotopy}. A formal equivalence induces a
quasi-isomorphism of $L$.

Continuous duality between discrete and linearly compact vector
spaces is exact. Since $L=(Q\mathcal B)^\vee[-1]$, this is equivalent
to (ii). The current recognition theorem identifies the recovered
map of $L$ with $1_{\mathfrak g}\otimes f_A$, proving (iii).

Finally filter cochains by $F^r$. The associated graded map is the
completed symmetric algebra on the map of linearized complexes.
Finite tensor products over $k$, finite symmetric coinvariants in
characteristic zero, and continuous duality preserve the relevant
quasi-isomorphisms. Every generator has positive cohomological degree,
so $F^{m+1}\mathcal B^m=0$. The filtration is bounded in each
degree, and filtered comparison proves the last assertion.
\end{proof}

\begin{example}\label{cycmp:ex:polynomial-collapse}
Ordinary field quasi-isomorphisms need not be formal equivalences.
Take $A=k[t]$, with $T_A=0$ and auxiliary weight
$\operatorname{wt}(t)=1$. Constant-current inclusion and evaluation
at $t=0$ give split field maps
\[
 \mathcal B_{k[t]}\rightleftarrows\mathcal B_k.
\]
Both are ordinary quasi-isomorphisms, including on every finite
weight quotient. To see this, fix chain weight $q$ and choose $N>q$.
The weight-$q$ Chevalley complex of $\mathfrak g[t]$ agrees with
that of $\mathfrak g[t]/t^N\mathfrak g[t]$: no bracket of total
weight at most $q$ meets the truncation.

For $\mathfrak g=\mathfrak{sl}_2$, the strong Macdonald theorem gives
exterior cohomology generators of degree three and cochain weights
\[
 0,-(N+1),\ldots,-(2N-1)
\]
\cite{cycmp:FGT}*{Theorem A and Remark 1.2(iii)}.
Thus every positive chain-weight component is acyclic, and weight
zero is exactly $\mathcal B_k$. Products of complexes of vector
spaces are exact, so the assertion also holds for the complete
cochains. Translation is zero, and the two maps preserve the scalar
map and the fields.

On the dual-number test, however, the first homotopy groups are
$\mathfrak g[t]\varepsilon$ and $\mathfrak g\varepsilon$.
Evaluation kills every positive polynomial weight. The coefficient
maps also fail to be equivalences: their singleton maps are inclusion
and evaluation between $\omega_X[1]$ and
$(\omega_X\otimes k[t])[1]$. Full coefficient duality reflects this
failure. Hence $\mathfrak F$ does not descend through ordinary
field quasi-isomorphisms, even with the scalar map retained.
\end{example}

Let $\mathcal P_{\mathfrak g}^{\mathrm{form}}$ denote localization
at formal equivalences, retaining all mapping spaces. Recognition
is an equivalence of relative categories, and therefore gives
\[
 \mathcal P_{\mathfrak g}^{\mathrm{form}}
 \simeq\bigl(\mathcal A[\mathrm{qis}^{-1}]\bigr)^{\mathrm{op}}.
\]
Direct products preserve quasi-isomorphisms, so the tensor structure
descends. These mapping spaces concern the specified quadratic
presentations with their split grading.

\begin{corollary}\label{cycmp:cor:formal-realization}
The coefficient realization descends to a covariant strong symmetric
monoidal functor on $\mathcal P_{\mathfrak g}^{\mathrm{form}}$.
It preserves both products, all supported collision comparisons,
every finite normal derivative, arbitrary spectators and every cut.
The field unit maps to $K$, the constant-current object maps to
$C_k^\sharp$, and the scalar maps and translation are preserved.
Full coefficient biduality recovers $C_A$ with these operations.
\end{corollary}

\begin{proof}
The formal equivalence theorem and the finite block filtration in
the proof of \eqref{cycmp:eq:colored-coefficients} show that
$\mathfrak F$ inverts the designated maps. Its structure comparisons
are natural before localization, so they descend with their
coherences and mapping spaces. Theorem~\ref{cycmp:thm:regular-realization}
identifies the products, units and translation.
The full retained coefficient duality of
Theorem~\ref{cyjets:thm:pro} returns every coefficient map on applying
the inverse duality.
\end{proof}

\begin{example}\label{cycmp:ex:contractible-jets}
Formal equivalence permits changes of presentation with all
translation jets. Set
\[
 J=k[x_0,x_1,\ldots]\otimes\Lambda(\eta_0,\eta_1,\ldots),
 \quad |x_j|=0,\quad |\eta_j|=-1,
 \quad d\eta_j=x_j,\quad dx_j=0,
\]
and define $Tx_j=x_{j+1}$, $T\eta_j=\eta_{j+1}$.
Both generators with index $j$ have weight $j+1$.
Every weight is a bounded finite-dimensional complex, so $J$ and
$A\otimes J$ belong to $\mathcal A$.

Let $N$ count total monomial degree, including the odd generators.
On positive monomial degree put
\[
 h=\frac1N\sum_{j\ge0}\eta_j\frac{\partial}{\partial x_j},
 \qquad h(1)=0.
\]
Every value is a finite sum. The supercommutator of
$d=\sum_jx_j\partial_{\eta_j}$ with
$\sum_j\eta_j\partial_{x_j}$ is $N$. Both preserve monomial degree,
so $dh+hd=1-\iota\epsilon$, where $\epsilon:J\to k$ is augmentation.

Translation commutes with $N$ and with both sums. Hence $Th=hT$.
The tensor homotopy is
$a\otimes v\mapsto(-1)^{|a|}a\otimes hv$.

Inclusion and augmentation between $A$ and $A\otimes J$ are therefore
quasi-isomorphisms commuting with translation. They give formal
equivalences of current fields and equivalences of the full
coefficient duals, with all scalar maps, vacua, cuts and collisions
preserved. The positive polynomial states in
Example~\ref{cycmp:ex:polynomial-collapse} admit no such contraction.
\end{example}

\section{Interacting currents and their rational coefficient space}
\label{cycmp:sec:interacting}

A quadratic current differential can split one current into two.
If the current carries an internal pole, that pole must survive the
split. We therefore keep the constituent jets distinct from their
grouping into currents. The differential changes the grouping and
leaves the constituent variables in place.

Work compatibly at every precision
$\mathsf R_r=k[\hbar]/(\hbar^r)$. Let $R$ be a finite-dimensional
commutative unital algebra, and let $E$ be a bounded nonpositive
complex of finite-dimensional $R$-modules. Put
\[
 F_R(E)=\operatorname{Sym}_R(E\otimes k[t]).
\]
The symmetric algebra uses the cohomological signs. A generator
$e_n=et^n$ has charge one and weight $n+1$. $R$ has charge and
weight zero.

Define even derivations, zero on $R$, by
\begin{equation}\label{cycmp:eq:jet-actions}
 T(e_n)=(n+1)e_{n+1},\qquad
 \delta_i(e_n)=e_{n-i+1},\qquad e_j=0\quad(j<0).
\end{equation}
They commute with the differential and obey
\[
 [\delta_i,\delta_j]=0,\qquad
 [\delta_i,T]=(i-1)\delta_{i-1},\qquad \delta_0=0.
\]
In particular, $\delta_1$ acts by charge.

\begin{definition}\label{cycmp:def:charged-quotient}
A charged jet quotient is an algebra $A=F_R(E)/\mathfrak a$, where
$\mathfrak a$ is a dg ideal of positive charge.
It is homogeneous in charge and weight and stable under $T$ and
every $\delta_i$.
Require each weight piece to be a bounded finite-dimensional complex.

Maps preserve the unit, differential, charge and jet actions.
The chosen presentation is not part of the data. Denote this category,
with the zero algebra admitted, by $\mathcal A_{\mathrm{cj}}$.
At precision $r$ use scalar extensions and $\mathsf R_r$-linear maps
preserving the same data.
\end{definition}

The algebra has a canonical presentation. On its charge-one part
$A_1$, put $D=\delta_2$ and $V_A=\ker D$. Since $[D,T]=1$ and
$D$ is locally nilpotent, the finite-input formulas
\begin{equation}\label{cycmp:eq:weyl-projector}
 \pi=\sum_{j\ge0}\frac{(-1)^j}{j!}T^jD^j,
 \qquad a=\sum_{n\ge0}\frac{T^n}{n!}\pi(D^na)
\end{equation}
give $A_1=V_A\otimes k[t]$. Adjacent binomial terms cancel to give
$D\pi=0$ and the second identity. Successive powers of $D$ prove
uniqueness.

Thus $V_A$ is a finite quotient complex of $E$, and
multiplication gives a canonical surjection $F_R(V_A)\to A$.
Its kernel is the exact relation ideal of this presentation.

Let $\mathcal C_A$ be the finite-weight polynomial core of the
Chevalley cochains of $\mathfrak g\otimes A$. On its generators put
\[
 T_B=-(1\otimes T)^\vee,\qquad
 D_i=-(1\otimes\delta_i)^\vee,
\]
and extend as derivations. These restricted transposes commute with
the current differential. Write $\mathcal C_A^{(n)}$ for total
charge $n$. On it $D_1=-n$.

Define
\begin{equation}\label{cycmp:eq:discriminant}
 \Lambda_n=k[x_1,\ldots,x_n]^{S_n},\qquad
 p_m=\sum_{a=1}^n x_a^m,
 \qquad d_n=\prod_{a<b}(x_a-x_b)^2.
\end{equation}
Set $\Lambda_0=k$ and $d_0=d_1=1$. The $x_a$ are constituent
dual-jet variables, independent of the positions on the curve.

\begin{lemma}\label{cycmp:lem:polarization}
On $\mathcal C_A^{(n)}$, the assignments $p_m\mapsto-D_{m+1}$
define a $\Lambda_n$-action. The current differential is linear
over this ring, multiplication by $d_n$ is injective, and
\[
 [T_B,p_m]=-mp_{m-1},\qquad T_B(d_n)=0.
\]
\end{lemma}

\begin{proof}
First take $A=F_R(E)$. Let $K_n$ be the multidegree
$(1,\ldots,1)$ current-cochain complex on independent copies
$E_1,\ldots,E_n$ of $E$. It is finite. A current uses a subset
of the copies, and a product of currents partitions them among its
factors. Constant-current ghosts are retained. Polarization gives
\begin{equation}\label{cycmp:eq:polarization}
 (K_n\otimes k[x_1,\ldots,x_n])^{S_n}
       \xrightarrow{\sim}\mathcal C_{F_R(E)}^{(n)}.
\end{equation}

Pair $x^j$ with $et^j$ by one, replace the copy indices of each
current by those jet indices, and use the graded symmetric pairing.
Repeated even indices contribute their factorials. Repeated identical
odd indices vanish. Each permutation orbit therefore maps bijectively
to its corresponding monomial in every finite weight. The
$R$-balancing relations concern the coefficient factors alone.

For example, if $R=E=k$, a cochain dual to the current
$a\otimes e_1\cdots e_n$ maps to $n!$ times the cochain dual to
$a\otimes e_0^n$. Multiplication under
\eqref{cycmp:eq:polarization} is concatenation followed by the
normalized projector
$P_{p+q}=((p+q)!)^{-1}\sum_{\sigma\in S_{p+q}}\sigma$.
There is no further binomial factor.

The quadratic differential splits a current's copy set into two
subsets. It changes no $x_a$. The linear differential acts on the
copies of $E$. In these coordinates
\[
 T_B=-\sum_a\partial_{x_a},\qquad
 D_{m+1}=-\sum_a x_a^m.
\]
The Newton identities, linearity and commutator formula follow.
The nonzero polynomial $d_n$ is injective before and after taking
invariants, and simultaneous translation annihilates it.

For a quotient, the state surjection induces an injection of current
cores. This follows on each finite weight and then on its symmetric
cochains. It preserves the differential and all jet
actions. The preceding identities and injectivity restrict to this
subcomplex. The moment formula makes the construction independent
of the chosen presentation.
\end{proof}

The meromorphic state complex is the intrinsic localization
\begin{equation}\label{cycmp:eq:meromorphic-core}
 \mathcal B_A^{\mathrm{mer}}=
 \bigoplus_{n\ge0}\mathcal C_A^{(n)}
                \otimes_{\Lambda_n}\Lambda_n[d_n^{-1}].
\end{equation}
Its differential extends uniquely. In a free presentation this is
$(K_n\otimes k[x_1,\ldots,x_n,\Delta_n^{-1}])^{S_n}$ in charge
$n$, where $\Delta_n=\prod_{a<b}(x_a-x_b)$: localizing at
$d_n=\Delta_n^2$ commutes with the averaging projector.
An internal pole therefore remains in the same coefficient space
when its current splits.

Choose
\[
 w(u)\in1+\hbar k[u,u^{-1}][[\hbar]]
\]
and reduce at each precision. For groups of charges $p,q$, set
\[
 \Omega_{p,q}(v)=\prod_{a=1}^p\prod_{b=1}^q w(v+x_a-y_b),
 \qquad R_{p,q}=\prod_{a,b}(x_a-y_b).
\]
These expressions are symmetric within the two groups.

Their joint
tensor is
\begin{equation}\label{cycmp:eq:joint-tensor}
 (\mathcal C_A^{(p)}\otimes\mathcal C_{A'}^{(q)})
                      [(d_pd_qR_{p,q}^2)^{-1}].
\end{equation}
Indeed the union of the two sets of variables sends $d_{p+q}$
exactly to $d_pd_qR_{p,q}^2$.

For field parameters also invert all
translated cross differences $x_a-y_b+mv$, $m\in\mathbb Z$.
With several parameters use their integral translates. The resulting
rational spaces are preserved by block translations and permutations.
Every element has finitely many denominator factors at fixed precision.

The ordinary current product $\mu_0$ is linear over
$\Lambda_{p+q}$, since each $D_i$ is primitive for multiplication.
It extends to \eqref{cycmp:eq:joint-tensor}. Define the native product
and fields by
\begin{equation}\label{cycmp:eq:meromorphic-product}
 \mu_w=\mu_0\Omega_{p,q}(0),\qquad
 Y_w(b,v)c=\mu_w(e^{vT_B}b\otimes c).
\end{equation}

\begin{proposition}\label{cycmp:prop:meromorphic-fields}
The product $\mu_w$ is unital, associative and differential graded.
Translation is locally nilpotent at each precision. The displayed
fields are regular dg nonlocal fields. On the joint rational space
their unpermuted exchange is
\begin{equation}\label{cycmp:eq:meromorphic-exchange}
 S_w(v)=\frac{\Omega_{q,p}(-v)^{21}}{\Omega_{p,q}(v)}.
\end{equation}
They satisfy both vacuum identities, translation, exchange,
unitarity, Yang--Baxter and the field hexagon. The construction is
contravariant in $\mathcal A_{\mathrm{cj}}$.
\end{proposition}

\begin{proof}
In a free presentation, the product is the normalized permutation
average of the ordinary product multiplied by $\Omega_{p,q}(0)$.
For three groups the final projector absorbs either intermediate
projector, and both associations have the factor
$\Omega_{p,q}(0)\Omega_{p,s}(0)\Omega_{q,s}(0)$.
The empty group gives the unit. The differential preserves the
constituent variables, so it commutes with every cross multiplier.
The ordinary dg product rule proves the new one.

For a quotient, the localized core embeds in the free localized
core. A two-group multiplier has numerator symmetric within each
group and denominator dividing a power of the total discriminant.
The numerator preserves the quotient-current factors through their
jet actions. Localized $\Lambda_{p+q}$-linear multiplication lands
in the quotient-current object, so all identities restrict to it.

Translation lowers polynomial weight and annihilates the internal
discriminant. It is locally nilpotent on every localized state.
Total translation annihilates the cross differences and is a
derivation of $\mu_w$. Associativity and this derivation give the
regular field identities, including both vacuum identities.

Keep the two groups unsymmetrized to check exchange. Translating
the first group by $-v$ changes $\Omega_{p,q}(v)$ to
$\Omega_{p,q}(0)$. The ratio in
\eqref{cycmp:eq:meromorphic-exchange} changes this factor to the
reversed product factor, while the final projector supplies the
Koszul permutation. Reversing the pair inverts the ratio.

The pair multipliers commute, proving Yang--Baxter. Splitting their
pair sets proves the hexagon at independent output separations.
Maps preserving the core jet actions preserve all these formulas
and extend uniquely through localization.
\end{proof}

At finite Taylor order, a translated denominator is inverted by its
finite nilpotent expansion. The compatible state tower has fields
in $\varprojlim_r\mathcal B_{A,r}^{\mathrm{mer}}[v]$.
The exchange space includes cross-group internal poles. Replacing
it by the ordinary Laurent tensor square would change the theorem.

\subsection{Recognition and the tensor of meromorphic currents}

A charged meromorphic current presentation retains its polynomial
quadratic current core, charge grading, scalar map, translation and
jet actions. Its completed core satisfies
Definition~\ref{cycmp:def:current-presentation}. On the recovered
states, charge zero is a finite algebra $R$ in degree zero.
Charge one has the decomposition \eqref{cycmp:eq:weyl-projector}
and generates the algebra over $R$. The relations satisfy
Definition~\ref{cycmp:def:charged-quotient}.

Its extension and native
operations are exactly \eqref{cycmp:eq:meromorphic-core} and
\eqref{cycmp:eq:meromorphic-product}. The core is marked data.
For polar $w$ it need not be a subalgebra for $\mu_w$.

Morphisms preserve core generators, differential, charge, scalar
map and jet actions, and extend by localization. Write this category
as $\mathcal G_w^{\mathrm{cj}}$. At finite precision use scalar
extensions and maps over $\mathsf R_r$.

\begin{proposition}\label{cycmp:prop:charged-recognition}
There is an equivalence
$\mathcal A_{\mathrm{cj}}^{\mathrm{op}}\simeq\mathcal G_w^{\mathrm{cj}}$.
It recovers the algebra, differential, unit, translation, every jet
action and the canonical relation ideal.
\end{proposition}

\begin{proof}
Current recognition recovers
$A_B=\operatorname{Hom}_{\mathfrak g}(\mathfrak g,W^\vee[-1])$.
A core derivation commuting with its differential and killing scalar
currents transposes to a current Lie derivation. Adjoint-isotypic
evaluation makes it $1\otimes\delta$ for a state derivation.
Thus negative transposition recovers all prescribed jet actions
and their commutators.

Charge zero recovers $R$, and \eqref{cycmp:eq:weyl-projector} recovers
$V_B=\ker(\delta_2|(A_B)_1)$. Generation gives the surjection
$F_R(V_B)\to A_B$ and its kernel. The symmetric-ring action and
localization are intrinsic, so they recover the entire meromorphic
object.

A morphism commuting with $\delta_2,T$ extends its map on
$V_B$ constantly along the jets. Its quadratic differential makes
it preserve multiplication and relations. Conversely every such
state map induces the core map and its unique localization. These
constructions are inverse on objects and morphisms.
\end{proof}

The tensor is also specified. On the joint localization write
$F_{B,C}=\Omega_{B,C}(0)$ and put
\[
 \sigma_w=\tau(F_{C,B}^{21})^{-1}F_{B,C},\qquad
 \mu_{B\otimes_w C}=(\mu_B\otimes\mu_C)
                           (1\otimes\sigma_w\otimes1),
\]
where $\tau$ is the Koszul permutation. If $U$ removes the native
cross factor, transport the polynomial core through
\begin{equation}\label{cycmp:eq:source-tensor}
 U_2:UB\otimes UC\longrightarrow U(B\otimes_w C),
 \qquad U_2=F_{B,C}^{-1}.
\end{equation}

\begin{lemma}\label{cycmp:lem:source-tensor}
This is a symmetric monoidal tensor corresponding to $A\times A'$.
Native multiplication is a morphism $B\otimes_w B\to B$, and
$U(\mu_w)U_2=\mu_0$. The tensor unit is
$\mathsf R_r=\mathcal B_{0,r}^{\mathrm{mer}}$.
\end{lemma}

\begin{proof}
The ordinary current core of $A\times A'$ is the tensor of their
cores. Its total discriminant gives exactly the joint localization.
For three groups,
\[
 F_{B,C\otimes D}=F_{B,C}F_{B,D},\qquad
 F_{B\otimes C,D}=F_{B,D}F_{C,D}.
\]
These equalities of commuting factors prove the associativity and
unit diagrams for $U_2$. Moving $\tau$ through the factors proves
the symmetry diagram and $\mu_w\sigma_w=\mu_w$.

For completeness, on four groups the braided product contributes
$w_{13}w_{24}w_{23}w_{32}^{-1}$. The factor on its two outputs
contributes $w_{12}w_{14}w_{32}w_{34}$. Their product is
$\prod_{i<j}w_{ij}$, the factor obtained by twisting within each
input and then between the inputs. Here $w_{ij}$ includes all
constituent pairs between groups $i,j$. This proves the algebra
tensor identity on the joint localization.

Finally $\mu_wF^{-1}=\mu_0$. The tensor constraint is the identity
on core generators, so the transported multiplication is the
recognized map corresponding to $A\to A\times A$.
\end{proof}

\subsection{The state fields and the full product comparison}

There is an independent field operation on the state algebra.
Put $\ell(u)=\log w(u)$ and define
\begin{equation}\label{cycmp:eq:state-contraction}
 \begin{split}
 C_\ell(u)&=\sum_{i,j\ge1}
  \frac{(-1)^{j-1}\ell^{(i+j-2)}(u)}{(i-1)!(j-1)!}
                      \delta_i\otimes\delta_j,\\
 \mathcal W_A(u)&=e^{C_\ell(u)},\qquad
 \mathcal E_A(u)=\mathcal W_A(u)^{-1}.
 \end{split}
\end{equation}
At precision $r$ the exponential has finitely many terms. On each
fixed state input only finitely many jet actions in each term survive.
The operators therefore have a finite pole bound on that input.

For $A\ne0$ set
\begin{equation}\label{cycmp:eq:state-fields}
 Y_A^w(a,u)b=\mu_A(e^{uT}\otimes1)\mathcal E_A(u)(a\otimes b),
 \qquad
 \mathcal S_A^w(u)=\mathcal E_A^{21}(-u)\mathcal E_A(u)^{-1}.
\end{equation}
Taylor translation is applied after contraction. It introduces
nonnegative powers and preserves the finite lower bound.

\begin{lemma}\label{cycmp:lem:state-fields}
These formulas define dg quantum vertex fields with vacuum $1_A$
and rational exchange at every finite precision. They are functorial
in $\mathcal A_{\mathrm{cj}}$.
\end{lemma}

\begin{proof}
The jet derivations commute. Their commutators with $T$ give
\[
 [T\otimes1,C_\ell(u)]=-\partial_uC_\ell(u),\qquad
 [T\otimes1+1\otimes T,C_\ell(u)]=0.
\]
This follows directly by shifting $i$ and $j$ in
\eqref{cycmp:eq:state-contraction}. The factorials cancel their
commutator coefficients. Every derivation annihilates the unit,
so a contraction with either vacuum is the identity.

The derivation rule splits a contraction with a product into its
two constituent contractions. Moving Taylor translation through it
shifts its separation by that parameter. Both three-field
associations have the same commuting pair factors before expansion.
Their reversed ratio is $\mathcal S_A^w$, proving associativity,
exchange and the field hexagon. Commutativity of pair factors proves
Yang--Baxter, and reversing a pair proves unitarity.

Each identity
is finite at fixed input, precision and Taylor coefficient.
The relation ideal is preserved by the derivations and multiplication,
so all operations descend and commute with the differential and maps.
\end{proof}

The native current multiplier determines this contraction. Undoing
translation in its first group changes $\Omega_{p,q}(0)$ to
$\Omega_{p,q}(u)$. With $u$ invertible its expansion is
\begin{equation}\label{cycmp:eq:moment-kernel}
 \log\Omega_{p,q}(u)=\sum_{a,b\ge0}
       \frac{(-1)^b\ell^{(a+b)}(u)}{a!b!}p_a(x)p_b(y).
\end{equation}
The charge terms $p_0(x)=p$, $p_0(y)=q$ are included.

Transposition of the moment actions recovers the two jet derivations.
Their two minus signs cancel. The finite-input transpose of
\eqref{cycmp:eq:moment-kernel} is exactly $C_\ell(u)$.

More explicitly, multiplication by $x_i$ transposes to lowering the
corresponding constituent jet index. If $N_{ij}$ is the difference
of the two lowerings, a factor $w(u+x_i-y_j)$ transposes to
\[
 w(u+N_{ij})a=\sum_{m\ge0}\frac{w^{(m)}(u)}{m!}N_{ij}^ma.
\]
The sum is finite on each jet input. Multiplying the factors gives
$\mathcal W_A(u)$, whose symmetric-ring expression descends through
the relations.

The normal-frame map retains density and derivatives:
\[
 [f(v)\,dv]\otimes a\longmapsto
 \sum_m[f^{(m)}(u)\,du]\otimes\frac{N_{ij}^ma}{m!}.
\]
It intertwines differentiation and lets $v$ act as $u+N_{ij}$.
No constituent variable is identified with a geometric position.

Let $C_A^0$ be the regular state collision complex and $C_A^w$ the
collision complex of \eqref{cycmp:eq:state-fields}. Both use the
partition coefficients, suspended signs and full normal modules
of the preceding sections. In a block order $\omega$, put
\[
 \mathcal E_A^\omega=
       \prod_{a<_\omega b}\mathcal E_{A,ab}(z_a-z_b).
\]
For ordered groups $G,H$, let $\mathcal E_{G,H}$ be the product
over all pairs with one input in each group, and put
$\mathcal W_{G,H}=\mathcal E_{G,H}^{-1}$.

\begin{lemma}\label{cycmp:lem:interacting-coefficients}
The ordered operators glue to an isomorphism
$\mathscr E_A:C_A^w\xrightarrow{\sim}C_A^0$.
For every ordered-subset decomposition $G\sqcup H$,
\begin{equation}\label{cycmp:eq:interacting-cut}
 (\mathscr E_G\boxtimes\mathscr E_H)\Delta^w_{G,H}
                \mathscr E_A^{-1}
             =\mathcal W_{G,H}\Delta^0_{G,H}.
\end{equation}
The cross operators commute with internal mergers, satisfy the
three-group coherence identities and preserve both local vacuum
actions, all normal derivatives and arbitrary spectators.
\end{lemma}

\begin{proof}
The translation commutators make each pair operator horizontal for
$\partial_{z_a}+T_a$, so it extends through the full closed direct
image. At an adjacent merger $z_a=z_b+u$, the derivation rule gives
\[
 \mu_{ab}e^{uT_a}
 \mathcal E_{ac}(z_b+u-z_c)\mathcal E_{bc}(z_b-z_c)
 =\mathcal E_{(ab)c}(z_b-z_c)\mu_{ab}e^{uT_a}.
\]
The corresponding identity holds in the other tensor place.
The internal pair factor is the contraction in
\eqref{cycmp:eq:state-fields}. Thus
$r^0_{ab}\mathcal E^\omega
=\mathcal E^{\omega'}r^w_{ab}$, before and after full principal
parts. No square-zero condition on positive-charge products is used.

Changing neighboring orders contributes
$\mathcal E_{ab}\mathcal E_{ba}^{-1}$, which is exactly the
quantum passage and cancels the change of ordered contractions.
Bring any nonconsecutive pair together to obtain the same merger
identity. The operators are even and preserve the suspended signs.

For an arbitrary subset, move $G$ before $H$ using these passages.
The order-change comparison cancels the passage, and the total
contraction factors as
$(\mathscr E_G\boxtimes\mathscr E_H)\mathcal E_{G,H}$.
This proves \eqref{cycmp:eq:interacting-cut}. Splitting pair sets
among three groups proves coherence, and the merger identity proves
compatibility with each internal differential. Horizontality retains
the density and every normal differential action. All contractions
with a vacuum are identities.
\end{proof}

For the outer merger $13$, write $u=z_1-z_3$, $t=z_2-z_3$ and
$X_{12}(u-t)=\mathcal E_{12}(u-t)\mathcal E_{21}(t-u)^{-1}$.
The identity includes the middle state explicitly:
\begin{equation}\label{cycmp:eq:interacting-outer}
 \begin{split}
 r^0_{13}\mathcal E_{12}(u-t)\mathcal E_{13}(u)\mathcal E_{23}(t)
 ={}&\mathcal E_{2,(13)}(t)\operatorname{pp}_{u=0}\mu_{13}e^{uT_1}\\
 &\hspace{6mm}\cdot\mathcal E_{13}(u)X_{12}(u-t).
 \end{split}
\end{equation}
It remains valid after multiplication by $u^{-m}$ for every $m\ge1$.
Further spectators supply the same pair factors on both sides.

Recover $A_B$ from the marked current core and set
$G(B)=(C_{A_B}^0)^\sharp$. Let $J^0$ be the transpose of the
colored-partition comparison combined with
\eqref{cycmp:eq:source-tensor}. On the $G\sqcup H$ component define
\begin{equation}\label{cycmp:eq:interacting-tensorator}
 J^w=J^0(\mathcal W_{G,H})^\sharp.
\end{equation}
The cross factor is fixed by the native kernel calculation
\eqref{cycmp:eq:moment-kernel}.

\begin{theorem}\label{cycmp:thm:interacting-realization}
The functor $G$ with tensorator $J^w$ is covariant and strong monoidal.
It has natural isomorphisms of full coefficient algebras
\[
 \mathscr E_A^\sharp:
 \mathfrak F_w(\mathcal B_A^{\mathrm{mer}},\mu_w)
          \xrightarrow{\sim}(C_A^w,\Delta^w)^\sharp.
\]
On every ordered-subset component the native product satisfies
\begin{equation}\label{cycmp:eq:interacting-product-square}
 G(\mu_w)J^w=(\Delta^0_{G,H})^\sharp
                                  (\mathcal W_{G,H})^\sharp.
\end{equation}
Both orders, all collisions and coherent cut compositions are retained.

The singleton is $(\mathscr M_A[1])^\sharp$. The field unit maps to
$K$, the constant-current object maps to $C_k^\sharp$, and $\rho_A$
maps to the local scalar projection. A chosen augmentation $R\to k$,
extended by zero on positive charge, gives its section.
Both local vacuum actions and their full normal modules are preserved.
\end{theorem}

\begin{proof}
Recognition and full coefficient duality each reverse arrows, making
$G$ covariant. The colored-partition comparison gives $J^0$, with
the source tensor of Lemma~\ref{cycmp:lem:source-tensor}.
For three colors both composites of the new tensorator contain
$\mathcal W_{G,H}\mathcal W_{G,J}\mathcal W_{H,J}$.
Empty groups contribute the identity. Lemma~\ref{cycmp:lem:interacting-coefficients}
makes these natural maps on every full open direct image, proving
the strong monoidal identities.

Untwisting native multiplication gives $\mu_0$ on the core. Its
generator transpose is the state diagonal $A\to A\times A$, whose
colored-partition map is $\Delta^0$. Hence
$G(\mu_w)J^0=(\Delta^0)^\sharp$. Inserting the cross factor proves
\eqref{cycmp:eq:interacting-product-square}. Transpose
\eqref{cycmp:eq:interacting-cut} to obtain the algebra isomorphism.

The merger identity, including \eqref{cycmp:eq:interacting-outer},
retains the full normal-transfer differential. Full coefficient
duality preserves the compositions and their coherent homotopies.

There is no pair contraction on a singleton. The zero state algebra
gives the empty unit, and $k\to A$ and the specified augmentations
preserve all jet actions. Their two transpositions give the stated
scalar maps. Contractions involving either scalar factor vanish,
which proves the vacuum assertions.
\end{proof}

Native fields are also actual morphisms after transporting their jet
marking. At finite Taylor precision give $B\{v\}$ the actions
\[
 D_i^{[v]}=e^{-vT_B}D_ie^{vT_B}
       =\sum_{j=0}^{i-1}\binom{i-1}{j}v^jD_{i-j}.
\]
Then $T_v=e^{vT_B}:B\{v\}\to B$ preserves the localized
presentation, and $Y_w(v)=\mu_w(T_v\otimes1)$.
State transposition recovers $a\mapsto e^{-vT}a$.

With $\Theta_v=C_A^0(e^{-vT})^\sharp$ and $m^w=G(\mu_w)J^w$,
monoidal naturality reads
\[
 \mathfrak F_w(Y_w(v))J^w_{B\{v\},B}
       =m^w(\Theta_v\star_\sharp1).
\]
The same conjugation changes $u$ to $u+v$ in the native kernel.
$v$ remains separate from all spatial coordinates.

The transported symmetry is
\[
 (J^w_{C,B})^{-1}G(\sigma_w)J^w_{B,C}
 =((\mathcal W_{H,G})^\sharp)^{-1}\tau_\sharp
                                      (\mathcal W_{G,H})^\sharp.
\]
Removing the permutation and inverting recovers the field convention
\eqref{cycmp:eq:meromorphic-exchange}. Reflection $w(u)\mapsto w(-u)$
gives opposite multiplication. Inversion $w\mapsto w^{-1}$ gives
inverse contraction. Full coefficient biduality returns $C_A^w$
with its density shifts, collisions and vacua. This is biduality
of the retained coefficient diagram.

\subsection{An interaction visible in cohomology}

Take
\[
 A_2=F_k(k)/(\text{charge at least three}),\qquad
 w(u)=1-\hbar/u,
\]
and write $a=e_0$. Its positive-charge ideal has nonzero square
and zero cube. For example,
\[
 Y_{A_2}^w(a,u)a=\frac{(e^{uT}a)a}{1-\hbar/u};
\]
the coefficient of $\hbar u^{-1}$ is $a^2$.
The polarized current complex retains partitions whose current
factors have at most two constituent indices. Let
$c_i^e,c_i^h,c_i^f$ be its odd charge-one current ghosts and put
\[
 \kappa_{12}=c_1^ec_2^f+c_1^fc_2^e+2c_1^hc_2^h,
 \qquad \beta=\frac{\kappa_{12}}{x_1-x_2},\qquad
 \gamma=(x_1+x_2)\beta.
\]
The ghost permutation sign and the odd denominator make these
degree-two states $S_2$-invariant.

\begin{proposition}\label{cycmp:prop:noncommutative-cohomology}
Over $\mathsf R_2$, the states $\beta,\gamma$ are closed and
$[\beta]*_w[\gamma]\ne[\gamma]*_w[\beta]$ in cohomology.
Thus this dg algebra is not quasi-isomorphic to a commutative dg
algebra over $\mathsf R_2$. The same holds for the free state
algebra $F_k(k)$ at this precision.
\end{proposition}

\begin{proof}
The invariant trace form gives $d\kappa_{12}=0$: its
constant-current action cancels, a charge-one current has no
positive-charge splitting, and the state differential vanishes.
The differential does not act on constituent variables, proving
closedness of both states.

On charge-four degree-four cochains, let $L$ extract the coefficient
of $c_1^hc_2^hc_3^hc_4^h$ and evaluate the variables at $(0,1,2,3)$.
All differences are invertible there. The functional kills every
boundary. A constant-current differential term retains a constant
ghost. A positive-current split could give the selected coefficient
only from a bracket of two Cartan currents, which is zero.

Set $f(x,y)=1/(x-y)$ and $g(x,y)=(x+y)/(x-y)$. The normalized
permutation average gives
\begin{equation}\label{cycmp:eq:cohomology-detector}
 \begin{split}
 L(\beta*_w\gamma-\gamma*_w\beta)
 ={}&-\frac{8\hbar}{4!}\sum_{\sigma\in S_4}\operatorname{sgn}(\sigma)
 f(x_{\sigma1},x_{\sigma2})g(x_{\sigma3},x_{\sigma4})\\
 &\quad\cdot\left(\sum_{a=1}^2\sum_{b=3}^4
              \frac1{x_{\sigma a}-x_{\sigma b}}\right)
          \bigg|_{(x_1,x_2,x_3,x_4)=(0,1,2,3)}
 =\frac{104}{9}\hbar.
 \end{split}
\end{equation}
The two Cartan coefficients contribute four, the difference of
first-order cross factors contributes minus two, and the projector
contributes $1/4!$. The nonzero value cannot be a boundary.
A commutative dg algebra has graded-commutative cohomology, so it
cannot be quasi-isomorphic to this algebra. The same detector and
boundary argument use only charge-one ghosts and remain valid for
the free state algebra.
\end{proof}

The formal criterion also extends to these marked currents. Complete
the polynomial core in weight, remove its native cross factor, and
test continuous maps to
$\mathsf R_r\otimes S\otimes\Omega_n$ carrying the core augmentation
ideal into $\mathsf R_r\otimes\mathfrak m_S\otimes\Omega_n$.
Here $S$ ranges over the nonpositive nilpotent test algebras used
in \eqref{cycmp:eq:formal-points}. Call a map a formal equivalence
when it induces weak equivalences on these tests for every $r,S$.

\begin{corollary}\label{cycmp:cor:interacting-formal}
A morphism of charged meromorphic current presentations is a formal
equivalence precisely when its recovered state map is a
quasi-isomorphism at every precision. The tensor, coefficient
realization and native field comparisons descend to localization
at these maps, retaining their mapping spaces and coherences.
\end{corollary}

\begin{proof}
The coordinate calculation in the formal equivalence theorem
identifies the tests with
$\operatorname{MC}_\bullet((\mathfrak g\otimes A_r)\otimes\mathfrak m_S)$.
The argument is over $k$, with the finite algebra $\mathsf R_r$
retained in its coefficients. Quasi-isomorphisms give equivalences
by nilpotence. Dual numbers have homotopy groups
$H^{1-n}(\mathfrak g\otimes A_r)\varepsilon$ for every $n\ge1$,
so they detect the converse.

The linear dual in this coordinate identification is exact.
The pairing $(a,b)\mapsto[\hbar^{r-1}](ab)$ identifies
$\mathsf R_r$ with its $k$-dual. In the power basis its matrix
is antidiagonal with entries one. Thus $\mathsf R_r$-linear duality
agrees with $k$-linear duality, weight by weight. The scalar-extended
finite weight pieces give the required continuous biduals.

The finite collision-block filtration sends a state
quasi-isomorphism to an equivalence of $C_A^0$, using derived
external products and geometric functors. Naturality of
$\mathscr E_A$ gives the same for $C_A^w$. Direct products preserve
state quasi-isomorphisms. Full duality and localization therefore
carry the tensor and the native field and product comparisons to
the localized categories. The universal property carries their
coherent compositions as well.
\end{proof}

\section{Current framings and the information they retain}
\label{cycmp:sec:framing}

The meromorphic construction retains its polynomial core, charge
and jet actions. Full coefficient duality preserves the geometric
current module, but a geometric current module need not select its
constant state sections. This distinction affects morphisms even
when every collision and vacuum map has been retained.

Before suspension the singleton is
$\mathscr M_A=\omega_X\otimes A$, with right action
\[
 (f\,dz\otimes a)\partial_z
       =-f'\,dz\otimes a-f\,dz\otimes Ta.
\]
An $\mathcal O_X$-linear map $F(z)$ is differential-linear exactly
when
\begin{equation}\label{cycmp:eq:horizontal}
 \partial_zF(z)+[T,F(z)]=0.
\end{equation}
A \emph{current framing} retains the constant state sections in
this presentation. A map induced by a translation-preserving state
map is constant in that framing. Equation \eqref{cycmp:eq:horizontal}
alone allows position-dependent maps.

\begin{proposition}\label{cycmp:prop:framed-reconstruction}
On the essential image of the preceding state-field construction,
retain a current framing and restrict arrows to actual constant
singleton maps. The binary boundary recovers the input fields by
\begin{equation}\label{cycmp:eq:decode}
 a_{(n)}b=[u^{-1}]r_Y(u^n;a,b),\qquad
 Y(a,u)b=\sum_{n\in\mathbb Z}(a_{(n)}b)u^{-n-1}.
\end{equation}
This recovery preserves translation, both vacuum identities and
the full collision and cut operations.

Suppose further that a comparison is supplied by an actual
isomorphism $f$ of recovered charged dg state algebras.
Require its singleton to be constant and horizontal and its state
map to preserve charge and scalar maps. It lifts
to a marked meromorphic current isomorphism if and only if it
preserves $\delta_2$.
\end{proposition}

\begin{proof}
Finiteness of each principal part gives the lower bound on field
powers. Function-linearity gives
\[
 [u^{-q-1}]r_Y(u^n;a,b)
       =[u^{-1}]r_Y(u^{n+q};a,b)\qquad(q\ge0).
\]
Thus decoding and re-encoding are inverse on every Laurent input.
Differential-linearity gives both translation identities with the
stated right-action sign. The two specified vacuum boundary
equations recover the two vacuum identities. A constant singleton
commutes with each coefficient extraction.

Full inverse duality retains every higher coefficient and normal
module. On an essential-image object, decoding returns its original
input field, so its full coefficient construction returns those
same collisions and cuts. The arrows here are induced by the stated
actual singleton maps, and their coefficient constructions are
natural. This recovers the original input fields. It does not
discard the marked core needed for the meromorphic realization.

Necessity of preserving $\delta_2$ follows from recognition.
For sufficiency, constant horizontality gives $[T,f]=0$.
On charge one, $\delta_i=\delta_2^{\,i-1}$, by
\eqref{cycmp:eq:weyl-projector}. Hence $f$ preserves every jet
there. The jet actions are derivations, vanish on the scalar
algebra, and charge one generates the algebra. The same identities
hold on all products and through the relation ideal.

The current transpose of $f$ preserves the split core and its
differential. It preserves the moment actions, so localization
extends it uniquely to each internal and joint rational space.
The native product, fields and tensorator then commute with it
term by term. Applying the construction to $f^{-1}$ gives the
inverse comparisons. This is a criterion for the supplied actual
state isomorphism, not a replacement for a derived lift by a test
on cohomology.
\end{proof}

\begin{example}\label{cycmp:ex:unmarked-translation}
Even a constant field automorphism can change the jet marking.
For the interacting algebra $A_2$ above, let $c\in\hbar\mathsf R_r$
and put $\psi_c=e^{cT}$. It is a finite exponential, preserves charge,
commutes with translation and has inverse $\psi_{-c}$.
Total translation commutes with every pair contraction. Thus
$\psi_c$ preserves the input fields and all their coefficient
operations.

Nevertheless
\[
 [\delta_2,\psi_c]=c\psi_c\delta_1,
 \qquad [\delta_2,\psi_c]e_0=c\psi_ce_0\ne0\quad(c\ne0).
\]
It is not an automorphism of the fixed marked presentation.

Transporting the marking gives
\[
 \delta_i^{(c)}=e^{cT}\delta_ie^{-cT}
     =\sum_{j=0}^{i-1}\binom{i-1}{j}(-c)^j\delta_{i-j}.
\]
With that transported marking $\psi_c$ is an isomorphism.
The assertion distinguishes a fixed marking from its isomorphic
transport. It does not distinguish the underlying unmarked field
algebra from itself.
\end{example}

\begin{proposition}\label{cycmp:prop:moving-framing}
For $A_2$ and every admissible $w$, the polynomial map
\[
 F_c(z)=\exp\bigl(c(z\delta_1+\delta_2)\bigr),
 \qquad c\in\hbar\mathsf R_r,
\]
induces a full coefficient automorphism, with inverse $F_{-c}(z)$.
It preserves the scalar maps, empty unit and both local vacua.
For $c\ne0$ it is not induced by a map of the fixed constant-state
fields preserving translation. Its fiber $f_c=e^{c\delta_2}$ at
$z=0$ satisfies
\begin{equation}\label{cycmp:eq:moving-framing}
 f_cTf_c^{-1}=T+c\delta_1,
 \qquad
 f_cY^w(f_c^{-1}a,u)f_c^{-1}b=e^{pcu}Y^w(a,u)b
\end{equation}
for an input $a$ of charge $p$.
\end{proposition}

\begin{proof}
The derivation $Z_2(z)=z\delta_1+\delta_2$ satisfies
$\partial_zZ_2+[T,Z_2]=0$. Its finite exponential is horizontal,
multiplicative and invertible. Explicitly,
\[
 F_c(z)e_n=e^{cz}\sum_{j=0}^n\frac{c^j}{j!}e_{n-j}.
\]
Every scalar exponential is finite in $\mathsf R_r[z]$.
The conjugation identity
$e^{uT}\delta_2e^{-uT}=\delta_2-u\delta_1$ gives
\[
 e^{uT}F_c(y+u)=F_c(y)e^{uT}.
\]

Each $F_c$ commutes with the jet derivations in both factors of
every contraction. Multiplicativity therefore gives the geometric
field identity
\begin{equation}\label{cycmp:eq:geometric-field}
 F_c(y)Y^w(a,u)b
       =Y^w(F_c(y+u)a,u)F_c(y)b.
\end{equation}
It also intertwines exchange.

Place $F_c(z_B)$ at each block of each partition coefficient.
Horizontal differential-linearity extends it through open and closed
direct images and every normal generator. Equation
\eqref{cycmp:eq:geometric-field} intertwines adjacent residues on
all Laurent multipliers. Exchange compatibility gives the same
identity for every nonconsecutive merger with spectators.

Cuts are external products of these same maps. All derivations
kill $1$ and preserve the augmentation ideal, proving the scalar
and vacuum assertions. The empty coefficient is fixed separately.
The negative exponential constructs the inverse. Reductions in
$r$ commute with these maps, giving a compatible formal tower.

A constant state map has position-independent singleton. But
$F_c(z)e_0=e^{cz}e_0$ has nonzero linear coefficient $ce_0$.
Its fiber also fails to commute with $T$:
$[f_c,T]=cf_c\delta_1$, since $[\delta_2,T]=\delta_1$ and the
next commutator vanishes. Conjugating the state-field formula
replaces $T$ by $T+c\delta_1$ and leaves its contraction unchanged.
This proves \eqref{cycmp:eq:moving-framing}.

At $r=2$, $c=\hbar$, the creation defect has coefficient
$\hbar e_0$ at $u^1$, and the interacting pair $(e_0,e_0)$ has
coefficient $\hbar e_0^2$. Both are nonzero in the degree-zero
state complex.
\end{proof}

Thus a reconstruction equivariant for all coefficient automorphisms
cannot recover these fixed constant-state fields with the supplied
framing. Full inverse duality still recovers the geometric current
contact system. The transported field with translation
$T+c\delta_1$ is isomorphic to the original through $f_c$. It is a
different translation-equipped presentation of that input.

The meromorphic current side detects the same distinction using its
own state complex. Full duality reverses compositions, so identify
a coalgebra automorphism $\Psi$ with the algebra automorphism
$(\Psi^{-1})^\sharp$. The inverse-transpose fiber of $F_c$ on
current generators is
\[
 g_c=(1\otimes f_c^{-1})^\vee=e^{cD_2}.
\]
It extends multiplicatively and through all joint localizations.
It preserves the native product, differential and scalar maps,
since $D_2$ commutes with the moment actions and is a product
derivation. But
\[
 g_cT_Bg_c^{-1}=T_B+cD_1,
 \qquad
 g_cY_w(g_c^{-1}b,v)g_c^{-1}d=e^{-qcv}Y_w(b,v)d
\]
for $b$ of charge $q$.

For the closed states $\beta,\gamma$ of the preceding section,
$D_2\beta=-\gamma$, $T_B\beta=0$, and $T_B\gamma=-2\beta$.
At precision two this gives
$g_\hbar\beta=\beta-\hbar\gamma$ and
$[g_\hbar,T_B]\beta=-2\hbar\beta$. The creation difference is
$-2\hbar v\beta$. Extracting $c_1^hc_2^h$ and evaluating
$(x_1,x_2)=(0,1)$ kills every boundary: a differential term
either contains a constant ghost or has Cartan coefficient
$[h,h]=0$. The functional takes $\beta$ to $-2$, and the
coefficient of $v$ in this defect to $4\hbar\ne0$.
The failure to preserve the fixed native field survives cohomology.

\begin{example}\label{cycmp:ex:joint-support}
The joint localization is also visible after a supported collision.
At precision two take $w(u)=1-\hbar/u$, $s=z_1-z_3$ and
$t=z_2-z_3$, with $t$ invertible. On the ordered cut
$G=\{1,3\}$, $H=\{2\}$, the merged state is $e_0^2$ at $z_3$
and the spectator is $e_0$ at $z_2$. Their separation is $-t$.
All higher jet actions vanish on these zero-jet states, so
\[
 \mathcal W_{G,H}(e_0^2\otimes e_0)
       =w(-t)^2(e_0^2\otimes e_0)
       =(1+2\hbar/t)(e_0^2\otimes e_0).
\]

The supported term lost by replacing this cross factor by the
identity is
\begin{equation}\label{cycmp:eq:joint-supported-term}
 \frac{2\hbar}{t}[ds/s]\otimes(e_0^2\otimes e_0).
\end{equation}
The two states lie in different tensor factors. Their tensor is
not the vanishing product $e_0^3$.

For $\eta_m=[ds/s^{m+1}]$, $m\ge0$, the full normal module has
\[
 \eta_m\partial_s=(m+1)\eta_{m+1},\qquad
 \eta_m s=\begin{cases}\eta_{m-1}&m>0,\\0&m=0.\end{cases}
\]
These actions give
$(\eta_m\partial_s)s-(\eta_ms)\partial_s=\eta_m$.
In particular $\eta_0$ and all its normal derivatives are retained.
The Laurent coefficient ring and the state modules are free over
$\mathsf R_2$, so \eqref{cycmp:eq:joint-supported-term} is nonzero.

Before the merger its cross factor is
$w(s-t)w(-t)=1-\hbar/(s-t)+\hbar/t$. On the normal test $s^{-m}$
its difference from the identity has principal part
\[
 \hbar\left(\frac{2}{t}s^{-m}
       +\sum_{j=1}^{m-1}\frac{s^{j-m}}{t^{j+1}}\right).
\]
Thus the comparison retains every normal order.

Already on two charge-one cores, the native tensor constraint is
$w(x-y)^{-1}=1+\hbar/(x-y)$. Their individual discriminants are
one, so the ordinary external tensor has no inverse of $x-y$.
The joint tensor inverts $(x-y)^2$ and contains the required term.
This domain difference precedes the supported coefficient difference
in \eqref{cycmp:eq:joint-supported-term}.
\end{example}

\section{Coherent descent and passage to a formal limit}
\label{cycmp:sec:descent}

The preceding comparisons give actual maps on full collision
diagrams. A global Hall-to-field comparison needs maps of this kind
on overlaps, with their product coherences and orientations. A
linear comparison of local complexes is only one part of that datum.
The cover must see every finite configuration. Its totalization
must retain the nonlinear equation governing the comparison.

A cover is a \emph{Weiss cover} when every finite subset is contained
in one of its members. An ordinary cover by two proper opens $U,V$
need not have this property: choose $p\in U\setminus V$ and
$q\in V\setminus U$. Their pair is contained in neither open.

Starting with a sufficiently fine polydisc basis subordinate to an
ordinary cover, take all finite disjoint unions of basis polydiscs.
For each finite subset choose pairwise disjoint small polydiscs
around its distinct points. Their union contains the subset, proving
the Weiss property. Such a union need not lie in a single member
of the original cover.

The Ran space records finite nonempty subsets, with its presentation
indexed by all nonempty finite sets and surjections. A surjection
$I\twoheadrightarrow J$ gives the diagonal $X^J\to X^I$.
There is no canonical map $\operatorname{Sym}^nX\to
\operatorname{Sym}^{n+1}X$ obtained by repeating a first entry:
the representatives $(p,q)$ and $(q,p)$ give different multiplicities.
The finite-set diagram retains the intended identifications.

Its configuration arity is independent of the dimension of a
\v Cech simplex. One insertion lying in three cover opens already
gives a two-simplex. A finite descent model therefore specifies
its normalized simplicial diagram and its coefficient complexes
separately from any arity or charge bound.

\subsection{The degree and product in the comparison equation}

Let $C$ be a coaugmented conilpotent dg coalgebra and $H$ a dg
associative algebra over a field of characteristic zero. Write
$\overline C=\operatorname{coker}(k\to C)$ and let
$\overline\Delta_C$ be its reduced coproduct. On
$\operatorname{Hom}(\overline C,H)$ use
\[
 Df=d_Hf-(-1)^{|f|}fd_C,
 \qquad f*g=\mu_H(f\otimes g)\overline\Delta_C.
\]
The tensor of maps uses the Koszul sign. The commutator of this
convolution product has degree zero.

\begin{lemma}\label{cycmp:lem:twisting-degree}
A degree-one map $t:\overline C\to H$ satisfying
$Dt+t*t=0$ determines a dg algebra map $\Omega C\to H$.
These assignments are inverse. For $C=BA$, this encodes a
comparison from the associative bar construction of $A$, with
its higher components.
\end{lemma}

\begin{proof}
Write $\Omega C=T(\overline C[-1])$. A generator $s^{-1}c$ has
degree $|c|+1$, so its image is $t(c)$. Take the cobar differential
on generators to be
\[
 d(s^{-1}c)=-s^{-1}(d_Cc)
       -\sum_{\overline\Delta c}
                  (-1)^{|c_1|}(s^{-1}c_1)(s^{-1}c_2).
\]
The coalgebra identities make its square zero.

The chain-map
equation on this generator is
\[
 d_Ht(c)+t(d_Cc)
       +\sum_{\overline\Delta c}(-1)^{|c_1|}t(c_1)t(c_2)=0,
\]
which is exactly $Dt+t*t=0$. The tensor algebra then extends the
map uniquely, and checking its generators suffices. Conilpotence
makes each reduced coproduct calculation finite.
\end{proof}

For a degree-zero map $\Theta$ in this same convolution algebra,
$[\Theta,\Theta]=0$. Its Maurer--Cartan expression cannot impose
multiplicativity. For instance the unital chain map on
$k[\epsilon]/(\epsilon^2)$ given by $f(\epsilon)=1$ is not an
algebra map. A shifted convolution convention is possible, but
requires the corresponding bracket, differential and map degrees.

On a finite semicosimplicial diagram of dg Lie algebras, retain
compatible polynomial forms on every simplex, as in the all-face
model of Chapter~\ref{cycuts:ch}. Differential and bracket are
componentwise. A compatible family of relative Maurer--Cartan
elements is then a Maurer--Cartan element of this total algebra.
Both the equation and its face restrictions are componentwise
equalities. Integration is a comparison of complexes. Transporting
the nonlinear equation to ordinary simplex cochains needs the
corresponding higher comparison maps.

\begin{example}\label{cycmp:ex:nonlinear-integration}
Linear integration need not preserve a Maurer--Cartan equation with
the Alexander--Whitney product. Let $V$ have an ordered basis of
degrees $(1,0,0,0)$ and zero differential. In $\operatorname{End}(V)$
put $A=E_{01}$ of degree one and $N=E_{12}+E_{23}$ of degree zero.
Then $N^3=0$, $NA=0$. Define
\[
 a(t)=A(1+tN+t^2N^2/2),\qquad
 \alpha=a(t)+N\,dt.
\]

The tensor differential gives $d\alpha=-a'(t)dt$, and
$\alpha^2=(a(t)N-Na(t))dt$. Since $a'=aN$ and $Na=0$,
$d\alpha+\alpha^2=0$.

Integration gives endpoint cochains
$a_0=A$, $a_1=A+AN+AN^2/2$ and edge cochain $N$.
The interval cochains have vertex idempotents $v_0,v_1$ and an
edge $e$ of degree one, with
\[
 dv_0=-e,\quad dv_1=e,\qquad v_0e=e=ev_1,
 \quad v_1e=ev_0=e^2=0.
\]
These define its Alexander--Whitney product.

In the signed tensor
product with the matrix algebra, the edge curvature is
\[
 -(a_1-a_0)+a_0N-Na_1=-AN^2/2=-E_{03}/2\ne0.
\]
Here the sign before the endpoint difference comes from the
degree-one matrix coefficient.

Replacing the edge by
$N+N^2/2=e^N-1$ makes the curvature zero and recovers
$a_1=a_0e^N$. The missing quadratic term is a nonlinear comparison
term. This calculation excludes the raw-integration inference.
It leaves the compatible-form construction intact.
\end{example}

Parameter dependence also affects the differential of a local
critical model. On $\mathbb A_x^1\times\mathbb A_t^1$ with
$W=tx$, the operator
\[
 D=\hbar d_x+d_xW\wedge+d_t
\]
satisfies $D^2(1)=-dx\wedge dt$. After inverting $\hbar$, the
operator
\[
 D'=\hbar d_x+d_t+d_xW\wedge
                      +\hbar^{-1}d_tW\wedge
\]
has square zero: the connection form is closed for
$\hbar d_x+d_t$, and its wedge square vanishes. The additional
term is $\hbar^{-1}x\,dt$, so it need not preserve a prescribed
$k[[\hbar]]$ lattice. A parameter-independent potential, a compatible
flat connection, or a specified larger model is therefore part of
the relative critical datum.

\subsection{Finite comparisons and inverse limits}

\begin{proposition}\label{cycmp:prop:finite-descent}
Let $f:C^\bullet\to D^\bullet$ be a map of semicosimplicial
complexes, with $C^p=D^p=0$ for $p>d$.
If each $f^p$ is a quasi-isomorphism, its finite totalization is a
quasi-isomorphism. If a sequence of such comparisons forms an
actual compatible tower $f_r:C_r\to D_r$, then it induces an
equivalence of derived inverse limits. When the tower maps are
degreewise surjective on both sides, ordinary inverse limits
compute these derived limits.
\end{proposition}

\begin{proof}
Filter the cone of the totalized map by semicosimplicial degree.
Its associated graded pieces are shifts of $\operatorname{Cone}(f^p)$.
There are finitely many, and each is acyclic, so the total cone is
acyclic.

For a tower of complexes $C_r$, a model for its derived limit is
\[
 \operatorname{Cone}\left(1-s:\prod_r C_r\longrightarrow
                                  \prod_r C_r\right)[-1],
\]
where $s$ uses the transition maps. Products of vector-space
complexes preserve quasi-isomorphisms, so the tower of comparisons
induces a quasi-isomorphism of these cones. If every transition map
is degreewise surjective, then $1-s$ is degreewise surjective:
solve its components successively, choosing a lift at each step.
Its kernel is the ordinary inverse limit. The short exact sequence
identifies that kernel with the displayed derived-limit model.
\end{proof}

For an operation-preserving comparison the operations must also be
compatible across this tower. In particular the tensor products,
normal-support completions and continuous maps on both sides must
be specified before passing to the limit. The proposition does not
identify a limit of ordinary tensors with a completed tensor.
Independent finite comparisons without compatible transition maps
do not define the tower to which it applies.

For a Hall-to-field application the finite inputs include oriented
critical complexes on extension correspondences and their actual
products. They also include local field complexes and comparison
maps on a common configuration diagram. Local chain equations, product equations
and higher overlap equations are separate conditions on those maps.

A cyclic contraction of finite models consists of maps
\[
 (M,d_M)\underset{p}{\stackrel{i}{\rightleftarrows}}(V,d),
 \qquad h:V\longrightarrow V[-1].
\]
Its underlying contraction identities are
\[
 pi=1,\qquad d i=i d_M,\qquad p d=d_Mp,
 \qquad ip+dh+hd=1.
\]
Its pairing compatibility and its compatibility with the extension
blocks are additional equations.

Finite bases turn supplied equations
into exact matrix equalities. They do not construct the contractions
or identify their coefficient spaces with vanishing cycles.

\begin{problem}\label{cycmp:qn:global-comparison}
Construct a comparison from independently oriented framed critical
Hall correspondences to an independently specified field recipient,
on a common full collision diagram. Require charge, the empty and
local units, all flag substitutions, both product orders, and the
supported outer coefficients to be preserved. Construct its coherent
overlap maps and compatible precision tower, with the stated
completion and normal transfer. Determine when the resulting
comparison is an equivalence and when a relative operation class
obstructs it under the admitted changes of orientation and comparison.
\end{problem}

The retained-exchange and marked-current constructions prove
reconstruction on their declared classes. The failures of raw
integration, an external tensor substitute, or an unframed field
recipient delimit those constructions. They do not decide the
question for every admissible field realization or every geometric
Hall comparison.
